% =============================================================================
% Chapter 2 - Background
% =============================================================================
% Included from main.tex via % =============================================================================
% Chapter 2 - Background
% =============================================================================
% Included from main.tex via % =============================================================================
% Chapter 2 - Background
% =============================================================================
% Included from main.tex via % =============================================================================
% Chapter 2 - Background
% =============================================================================
% Included from main.tex via \include{contents/chapter2_shortened}.
% This file is the single source for Chapter 2; the older chapter2.tex and the
% chapter2_short_half*.tex files are superseded and must not be included.
%
% Equations are numbered and labelled eq:*; each is followed by a definition of
% its symbols. The mathematics of every equation is unchanged from the version
% reproduced in contents/formula_guide_complete.tex.
%
% Figure 2.1 (fig2_1_casme2_sequence.png) is Figure 2 of Yan et al. (2014),
% PLoS ONE 9(1):e86041, reproduced under the CC BY licence with attribution.
%
% Requires in the preamble: amsmath, graphicx.
% The pipeline overview (formerly the last section) now lives in Chapter 4.
% =============================================================================

\chapter{Background}
\label{ch:background}

This chapter presents the concepts on which the remainder of the thesis builds. \autoref{sec:phenomenon} defines micro-expressions and their characteristics. \autoref{sec:recording-to-input} describes the recordings of the CASME~II dataset and their preparation as model input. \autoref{sec:motion-magnification} introduces Eulerian video magnification, and \autoref{sec:video-to-motion} the optical flow and optical strain representations derived from the magnified video. \autoref{sec:nn-fundamentals} and \autoref{sec:building-blocks} cover the neural network concepts and building blocks used by the model. \autoref{sec:scarcity-skew} addresses training with limited and imbalanced data, and \autoref{sec:evaluation} presents the evaluation methodology.


\section{Micro-Expressions}
\label{sec:phenomenon}

Facial expressions are commonly divided into macro-expressions and micro-expressions. Macro-expressions are the ordinary expressions of everyday interaction; they typically last between 0.5 and 4 seconds and involve a large area of the face~\cite{ben2021,xie2022,adegun2020}. Micro-expressions are brief, involuntary and local facial movements that tend to occur when a person attempts to conceal or suppress an emotion, particularly in high-stakes situations~\cite{yan2014,ben2021,xie2022,li2022survey}. Ekman and Friesen described such movements as a form of nonverbal leakage~\cite{ekman1969}, and they have since been studied empirically in concealed and falsified emotional expressions~\cite{porter2008}. Their involuntary character distinguishes them from expressions deliberately posed for a camera~\cite{qu2016}. Comprehensive reviews of automatic micro-expression analysis are provided by Oh et al.~\cite{oh2018}, Ben et al.~\cite{ben2021}, Li et al.~\cite{li2022survey} and Xie et al.~\cite{xie2022}.

\autoref{fig:casme2-sequence} shows an example from CASME~II: a lowering of the inner brow (AU~4, together with AU~9) labelled as disgust, whose apex occurs approximately 110~ms after onset~\cite{yan2014}. The change between frames is barely visible at full-face scale, which illustrates why the phenomenon is difficult to perceive. An observed movement is nevertheless neither proof of deception nor an unambiguous indicator of emotion; the recognition task depends on the annotations of the dataset.

\begin{figure}[htbp]
	\centering
	\includegraphics[width=\textwidth]{fig2_1_casme2_sequence.png}
	\caption[Frame sequence of a CASME II micro-expression]{Frame sequence of a micro-expression from CASME~II, with time in milliseconds from onset. The action units are AU~4 and AU~9, labelled as disgust; the apex occurs at approximately 110~ms. The insets magnify the right inner brow (AU~4). Reproduced from Yan et al.~\cite{yan2014} (Figure~2) under the Creative Commons Attribution (CC~BY) licence.}
	\label{fig:casme2-sequence}
\end{figure}

\subsection{Duration}
\label{sec:duration}

Definitions of micro-expression duration vary. Yan et al.~\cite{yan2014} describe half a second as the generally accepted upper limit, whereas other sources report typical durations of 1/25 to 1/3 of a second~\cite{li2022survey} or 0.065 to 0.5 seconds~\cite{ben2021}. The limit is a working definition rather than a duration filter applied by this pipeline. A brief event yields few observations at ordinary video frame rates, so capture frequency and temporal sampling determine the evidence available for recognition.

\subsection{Intensity}
\label{sec:low-intensity}

Suppressed facial movements may be difficult to perceive even when sufficient frames are recorded~\cite{yan2014}. Because micro-expressions are local~\cite{ben2021}, a small contraction may affect only the region around the brow or mouth while most of the image remains unchanged. Duration limits the number of observations, whereas low intensity limits how clearly the movement can be separated from noise within them. Both properties motivate motion magnification and explicit motion features.

\subsection{Temporal Phases: Onset, Apex and Offset}
\label{sec:onset-apex-offset}

The \textbf{onset} marks the beginning of the movement, the \textbf{apex} its point of greatest annotated intensity, and the \textbf{offset} its end (\autoref{fig:casme2-sequence}). Apex-frame methods use the peak frame as a compact description~\cite{liY2018}. This pipeline instead uses the annotated onset and offset to bound a sequence. The apex annotation is neither read by the model nor used to select frames, so temporal features must be learned from the sampled sequence rather than supplied through a selected peak.

\subsection{Action Units and the Prediction Target}
\label{sec:action-units}

The Facial Action Coding System (FACS)~\cite{ekman1978} describes facial movements through Action Units (AUs), each corresponding to an individual component of muscle movement~\cite{li2022survey}, rather than through emotion categories. CASME~II assigns emotion labels using coded movements together with other annotation evidence~\cite{yan2014}, and the mapping is not one-to-one. This thesis predicts the resulting emotion category only. AU metadata is retained in the coding table but is neither an input nor a training target, and the system does not perform AU recognition.


\section{Data Acquisition and Preprocessing}
\label{sec:recording-to-input}

Input preparation determines which recorded information reaches every subsequent component. In particular, image resizing must be distinguished from face alignment and from temporal interpolation.

\subsection{High-Speed Recording}
\label{sec:high-speed-capture}

The number of frame intervals spanned by a facial movement is approximately
\begin{equation}
	n \approx D f_s,
	\label{eq:frame-intervals}
\end{equation}
where $n$ is the number of frame intervals, $D$ is the duration of the movement in seconds and $f_s$ is the recording rate in frames per second (fps). For an event lasting $D = 0.3$~s, this gives roughly eight to nine intervals at 25--30~fps, against sixty at 200~fps. CASME~II is recorded at 200~fps~\cite{yan2014}. The pipeline takes this rate from its configuration rather than reading it from each clip. After temporal subsampling, the physical spacing between selected frames need not match the nominal rate.

\subsection{Source Frames}
\label{sec:corpus-supplies}

CASME~II provides both the original recorded frames and a registered, cropped release. The study manifest identifies folders of the original $640\times480$ frames rather than the normalised release. The preprocessing performs no face detection, landmark localisation, registration, alignment or cropping of its own. The input therefore includes the full recorded field of view, and neither facial position nor scale is geometrically standardised. The registration described by Yan et al.~\cite{yan2014} must not be attributed to this pipeline.

\subsection{Frame Preparation}
\label{sec:frame-preparation}

Frames are loaded in greyscale and resized by bilinear interpolation to $224\times224$ pixels. Resizing the whole rectangular frame to a square changes its aspect ratio; the result is not a cropped face region. The loader selects the frames between the annotated onset and offset and then samples a fixed-length sequence. No apex localisation is performed.

\subsection{Clip Selection}
\label{sec:clip-selection}

Preprocessing selects CASME~II entries that are coded as micro-expressions, marked as having frames, and have an inclusive onset-to-offset length greater than two frames. The loader must still find readable frames, since the metadata flag alone does not guarantee successful extraction. The recognition dataset additionally requires a matching saved tensor and a label in the selected three-class mapping; the residual \texttt{Others} category is excluded. Preprocessing eligibility and final classification membership are therefore separate filters.


\section{Eulerian Video Magnification}
\label{sec:motion-magnification}

\subsection{The Eulerian Approach}
\label{sec:eulerian-idea}

Eulerian video magnification (EVM), introduced by Wu et al.~\cite{wu2012}, examines intensity changes at fixed image locations rather than explicitly tracking moving points. At each pixel, it isolates temporal variation within a selected frequency band, scales that variation and adds it back to the signal. Bai et al.~\cite{bai2021} describe the amplitude-based mechanism through the small-motion approximation
\begin{equation}
	I(x,t)=f(x+\delta(t))
	\approx f(x)+\delta(t)\frac{\partial f}{\partial x},
	\label{eq:evm-small-motion}
\end{equation}
where $I(x,t)$ is the image intensity at position $x$ and time $t$, $f$ is the intensity profile of the image at the initial time, and $\delta(t)$ is the displacement of that profile at time $t$. Adding $\alpha$ times the motion-related variation gives approximately
\begin{equation}
	\widetilde I(x,t)\approx
	f(x)+(1+\alpha)\delta(t)\frac{\partial f}{\partial x},
	\label{eq:evm-magnified}
\end{equation}
where $\widetilde I(x,t)$ is the magnified intensity and $\alpha$ is the amplification factor. For sufficiently small displacements, intensity amplification therefore resembles a displacement enlarged by the factor $(1+\alpha)$. The approximation does not imply that arbitrary intensity changes represent facial movement, nor that larger amplification always improves recognition.

\subsection{Spatial Decomposition with the Laplacian Pyramid}
\label{sec:laplacian-pyramid}

A Laplacian pyramid~\cite{burt1983} separates an image into spatial scales. At each level, the image is blurred and downsampled, and an upsampled reconstruction is subtracted from the current image to retain the detail at that scale. Repeating this procedure produces a set of Laplacian detail bands and a final coarse Gaussian residual. Reconstruction starts from the residual, upsamples it and adds the detail bands successively. The implementation uses the separable blur kernel $[1,4,6,4,1]/16$ and factor-two downsampling.

\subsection{Temporal Filtering}
\label{sec:temporal-filtering}

Each pixel in a spatial band supplies a time series. A band-pass filter retains the variation within the selected temporal-frequency interval, and the retained signal is amplified before reconstruction. Bai et al.~\cite{bai2021} describe a Butterworth filter for amplitude-based magnification. The present implementation instead sets to zero the Fourier coefficients outside the selected band and then applies the inverse transform.

\subsection{The Amplification Trade-Off}
\label{sec:amplification-tradeoff}

The filter distinguishes frequencies, not causes. Head movement, illumination variation and noise within the pass-band may be amplified together with the facial contraction of interest. A larger $\alpha$ can also violate the small-displacement approximation of \eqref{eq:evm-small-motion} and produce artefacts. Magnification therefore changes the signal presented to flow estimation without guaranteeing a cleaner or more discriminative representation.

\subsection{Implementation Details}
\label{sec:evm-departures}

\autoref{tab:ch2-evm-departures} summarises the implementation choices that affect the interpretation of the magnified input.

\begin{table}[htbp]
    \centering
    \small
    \begin{tabular}{@{}p{0.26\linewidth} p{0.68\linewidth}@{}}
        \hline
        \textbf{Property} & \textbf{Implemented behaviour and consequence} \\
        \hline
        Temporal filter & Hard Fourier mask rather than the Butterworth filter described by Bai et al.~\cite{bai2021}; the pass-band has no smooth roll-off. \\
        Spatial amplification & The same $\alpha$ is used for every Laplacian band, without spatial-frequency damping. The Gaussian residual is neither filtered nor amplified. \\
        Colour treatment & Processing operates on greyscale intensity, with no separate chromatic attenuation. \\
        Pipeline ordering & Each clip is sampled to 33 frames before magnification, but the filter still assumes the nominal 200 fps. Its physical pass-band is therefore clip-dependent. \\
        \hline
    \end{tabular}
    \caption[EVM implementation departures]{Implementation choices that distinguish the magnifier and its position in the pipeline from the reference formulation. These choices are shared by all magnification-enabled configurations.}
    \label{tab:ch2-evm-departures}
\end{table}

The Laplacian pyramid itself is retained. After reconstruction, values are clipped to $[0,255]$ and converted to 8-bit greyscale before motion extraction. Because subsampling changes the frame spacing, the nominal band does not correspond to the same physical-frequency interval for every clip.


\section{Motion Representation: Optical Flow and Optical Strain}
\label{sec:video-to-motion}

\subsection{Motivation for Motion-Based Input}
\label{sec:why-motion}

Optical flow makes displacement explicit instead of requiring the network to infer it from appearance alone. This can reduce reliance on raw image intensity but does not guarantee invariance to identity or illumination: facial geometry, registration errors and brightness changes still affect the estimated field. The resulting representation therefore remains conditioned on the capture and preprocessing procedures.

\subsection{Optical Flow and Brightness Constancy}
\label{sec:optical-flow}

Optical flow estimates the apparent displacement of image content between frames~\cite{horn1981}. It rests on the brightness-constancy assumption that a moving point retains its intensity, as described for micro-expression recognition in OFF-ApexNet~\cite{liong2019a}:
\begin{equation}
	I(x,y,t)=I(x+u\Delta t,y+v\Delta t,t+\Delta t),
	\label{eq:brightness-constancy}
\end{equation}
where $I(x,y,t)$ is the intensity at pixel $(x,y)$ and time $t$, $u$ and $v$ are the horizontal and vertical velocity components, and $\Delta t$ is the time interval between frames. A first-order Taylor expansion gives the optical flow constraint equation~\cite{horn1981}
\begin{equation}
	I_xu+I_yv+I_t=0,
	\label{eq:flow-constraint}
\end{equation}
where $I_x$, $I_y$ and $I_t$ are the partial derivatives of the intensity with respect to $x$, $y$ and $t$. This is a single equation in two unknowns. Under the \textbf{aperture problem}, a local edge constrains the motion normal to the edge, along the intensity gradient, but does not determine the motion parallel to it. Estimators therefore introduce neighbourhood assumptions or regularisation. In the discrete pipeline, $u$ and $v$ denote displacement in pixels between selected frames rather than physical velocity per second. Illumination changes and large inter-frame displacements can violate the underlying approximation.

\subsection{Optical Strain}
\label{sec:optical-strain}

Optical strain describes how neighbouring displacements differ. For the displacement field $\mathbf{u}=[u,v]^{\mathsf T}$, the small-strain approximation used for optical strain~\cite{shreve2011} is
\begin{equation}
	\varepsilon=\tfrac12\left(\nabla\mathbf{u}
	+(\nabla\mathbf{u})^{\mathsf T}\right),
	\qquad
	\varepsilon_{xx}=\frac{\partial u}{\partial x},\quad
	\varepsilon_{yy}=\frac{\partial v}{\partial y},
	\label{eq:strain-tensor}
\end{equation}
\begin{equation}
	\varepsilon_{xy}=\varepsilon_{yx}
	=\tfrac12\left(\frac{\partial u}{\partial y}
	+\frac{\partial v}{\partial x}\right),
	\label{eq:strain-shear}
\end{equation}
where $\varepsilon$ is the $2\times2$ strain tensor, $\nabla\mathbf{u}$ is the spatial gradient (Jacobian) of the displacement field, $\varepsilon_{xx}$ and $\varepsilon_{yy}$ are the normal strain components in the horizontal and vertical directions, and $\varepsilon_{xy}=\varepsilon_{yx}$ is the shear strain component. The normal components describe local stretching or compression, and the off-diagonal terms describe shear. A spatially uniform translation has zero gradient and therefore zero strain. This property does not imply immunity to arbitrary head rotation, projection effects or flow-estimation errors, and differentiation may amplify noise.

The implemented scalar strain magnitude is
\begin{equation}
	\varepsilon_{\mathrm{mag}}
	=\sqrt{\varepsilon_{xx}^2+\varepsilon_{yy}^2
		+\varepsilon_{xy}^2},
	\label{eq:strain-magnitude}
\end{equation}
where $\varepsilon_{\mathrm{mag}}$ is the strain magnitude at a pixel and the remaining symbols are as defined in \eqref{eq:strain-tensor} and \eqref{eq:strain-shear}. This form counts shear once. The Frobenius norm of the symmetric tensor counts both off-diagonal entries, giving $2\varepsilon_{xy}^2$ inside the root. The two forms differ in the relative weighting of shear, not merely by a removable global scale. The scalar discards the direction of deformation but accompanies the original displacement components as a third channel.

\subsection{Construction of the Three-Channel Input Tensor}
\label{sec:three-channel-tensor}

Uniform index sampling selects 33 existing frames from the onset-to-offset sequence. No Temporal Interpolation Model (TIM)~\cite{pfister2011} or other frame synthesis is performed, and repeated indices can occur for short clips. After optional magnification, Farneb\"ack optical flow~\cite{farneback2003} is computed for the 32 adjacent frame pairs, following the estimator choice reported by Zhao et al.~\cite{zhaoS2021}. The OpenCV settings are pyramid scale 0.5, three levels, window size 15, three iterations, polynomial neighbourhood 5, polynomial smoothing 1.2 and flags 0.

The spatial derivatives required for strain are computed with NumPy's \texttt{gradient} function, using unit pixel spacing, central differences in the interior and first-order one-sided differences at the boundaries. Temporal derivatives are not used, and no dedicated strain post-filter is applied. Stacking $[u,v,\varepsilon_{\mathrm{mag}}]$ gives a tensor of shape $(3,32,224,224)$ over the resized full frame. The network receives this tensor rather than the original intensity sequence.

Two normalisations are applied. During extraction, each channel is mapped independently to $[0,1]$ using its extrema across the clip. During loading, each channel is standardised over its $T\times H\times W$ values, after any training augmentation. Small constants in the denominators protect nearly constant channels, which therefore need not have unit variance. Both normalisations operate per sample and use no statistics pooled across training and held-out subjects.


\section{Neural Network Fundamentals}
\label{sec:nn-fundamentals}

A neural network composes parameterised operations to map the motion tensor to emotion-class scores~\cite{lecun2015,goodfellow2016}. This section introduces only the concepts required to interpret the implemented blocks and their controls.

\subsection{Units, Layers and the Forward Pass}
\label{sec:units-layers}

A unit computes
\begin{equation}
	a=\phi(w^{\mathsf T}x+b),
	\label{eq:unit}
\end{equation}
where $x$ is the input vector, $w$ is the learned weight vector, $b$ is the learned bias, $\phi$ is an activation function and $a$ is the resulting activation~\cite{goodfellow2016}. A layer applies several such operations in parallel, and composing layers gives a network. Depth refers to the number of layers and width to their feature dimension. A forward pass evaluates these operations with the current parameters and produces class scores, or logits, from which class probabilities are obtained.

\subsection{Non-Linear Activation Functions}
\label{sec:non-linearity}

Without non-linear operations, a stack of affine layers reduces to a single affine map. Non-linearity allows the composition to represent more complex relationships. The convolutional streams use the rectified linear unit (ReLU)~\cite{nair2010}, and the transformer feed-forward layers use the Gaussian error linear unit (GELU)~\cite{hendrycks2016}:
\begin{equation}
	\operatorname{ReLU}(x)=\max(0,x),
	\qquad
	\operatorname{GELU}(x)=x\,\Phi(x),
	\label{eq:activations}
\end{equation}
where $x$ is the input to the activation and $\Phi$ is the cumulative distribution function of the standard normal distribution. Both are fixed functions, although they act on features learned by the surrounding layers.

\subsection{Gradient-Based Learning}
\label{sec:gradient-descent}

A loss function compares predictions with labels. Backpropagation~\cite{rumelhart1986} applies the chain rule through the network to compute the derivatives of the loss with respect to its parameters. An optimiser uses these gradients to update the parameters, with the learning rate controlling the size of each update. Mini-batches provide successive updates, and an epoch completes the configured number of batch draws. When sampling is performed with replacement, an epoch need not visit every distinct clip.

During evaluation, the parameters are held fixed and training-only stochastic behaviour such as dropout is disabled. A low training loss alone does not establish useful learning, because the relevant objective is prediction on unseen subjects. A validation set used to select an epoch therefore has a different role from a test set reserved for final assessment.

\subsection{Convolution}
\label{sec:convolution}

A convolution applies the same learned kernel at multiple spatial positions~\cite{lecun1998}. This weight sharing allows local patterns to be detected without learning a separate detector for every image location. Each output channel is a feature map, and each kernel combines all input channels over its local window. For a two-dimensional layer with one bias per output channel, the number of parameters is
\begin{equation}
	N_{\mathrm{params}}=C_{\mathrm{out}}
	(C_{\mathrm{in}}k_Hk_W+1),
	\label{eq:conv-params}
\end{equation}
where $C_{\mathrm{in}}$ and $C_{\mathrm{out}}$ are the numbers of input and output channels, $k_H$ and $k_W$ are the kernel height and width, and the added one accounts for the bias of each output channel. The parameter count is independent of the image dimensions, but computation and activation memory increase with the number of positions processed. The stride sets the spacing between kernel applications, and padding controls boundary treatment and output size. Stacking layers enlarges the receptive field, that is, the input region that can influence a later activation.

\subsection{Pooling}
\label{sec:pooling}

Pooling replaces a neighbourhood with a fixed summary such as its maximum or mean~\cite{lecun1998,goodfellow2016}. A $2\times2$ max-pooling operation with stride two halves the spatial dimensions and retains the strongest response in each window. Adaptive average pooling instead produces a specified output grid. Neither operation learns weights. Both discard detail while reducing the computation of subsequent layers; any later learned projection is a separate operation with its own parameters.

\subsection{Model Capacity and Small Datasets}
\label{sec:capacity}

Capacity describes the range of functions a model can represent~\cite{goodfellow2016}. The parameter count is one indicator, but architecture, constraints and training procedure also matter. A model can overfit subject-specific appearance, noise or recording conditions, achieving a low training loss without generalising. Having more parameters than examples does not by itself prove overfitting, just as few parameters do not guarantee good predictions.

The small size of the CASME~II training pool motivates controlled architectural comparisons and regularisation. Their effectiveness must be assessed on held-out subjects, with checkpoint selection kept separate from final testing where possible. Because the analytical flow and strain transformations already expose motion structure, the learned network does not start from raw pixels.


\section{Network Building Blocks}
\label{sec:building-blocks}

The model combines a spatial stem, optional activation re-weighting, an optional temporal encoder and a common classification head. The stem consists of three streams with separate learned weights, one each for horizontal flow, vertical flow and strain. Disabling a component replaces a specified operation while preserving compatible input and output shapes.

\subsection{Three-Dimensional Convolution and Kernel Shape}
\label{sec:conv3d}

A three-dimensional convolution~\cite{ji2013,tran2015} operates on tensors of shape $(B,C,T,H,W)$, denoting batch size, channels, time, height and width. A kernel of shape $(k_T,k_H,k_W)$ mixes information across time only when $k_T>1$. The implemented streams use $(1,3,3)$ kernels with $(0,1,1)$ padding, followed by $(1,2,2)$ max-pooling. The convolutions therefore preserve the temporal length and learn spatial filters, while pooling halves the height and width.

The use of \texttt{Conv3d} does not establish that temporal convolution was tested. The complete stream can nevertheless depend on multiple frames, because batch normalisation during training and SimAM, whenever enabled, use statistics that span time. This distinction separates temporal convolution from other forms of whole-clip dependence.

\subsection{Normalisation and Regularisation}
\label{sec:normalisation}

Batch normalisation~\cite{ioffe2015} rescales each channel using a mean and variance, followed by a learned scale and shift:
\begin{equation}
	\widehat x=\frac{x-\mu_c}{\sqrt{\sigma_c^2+\epsilon}},
	\qquad y=\gamma_c\widehat x+\beta_c,
	\label{eq:batchnorm}
\end{equation}
where $x$ is an activation in channel $c$, $\mu_c$ and $\sigma_c^2$ are the mean and variance of that channel, $\epsilon$ is a small constant for numerical stability, $\widehat x$ is the normalised activation, $\gamma_c$ and $\beta_c$ are the learned scale and shift of the channel, and $y$ is the output. \texttt{BatchNorm3d} computes the training statistics over the $(B,T,H,W)$ dimensions of each channel, and evaluation uses running estimates accumulated during training. These activation statistics differ from the per-sample input standardisation performed by the data loader.

Dropout~\cite{srivastava2014} randomly deactivates units during training to reduce overfitting. \texttt{Dropout3d} removes entire feature channels for a sample~\cite{tompson2015} rather than masking individual spatial values. The retained channels are rescaled, and evaluation disables the random mask. Layer normalisation~\cite{ba2016} instead standardises the feature dimension of each token independently of other samples. It appears inside the transformer and in the classification head, where it is followed by dropout and a linear projection to class logits.

\subsection{Parameter-Free Attention: SimAM}
\label{sec:simam}

SimAM assigns larger weights to activations that differ from their channel context~\cite{yang2021}. Following the reference code of the original paper, the statistics of each context are
\begin{equation}
	\mu=\frac1M\sum_i x_i,\qquad
	v=\frac1{M-1}\sum_i(x_i-\mu)^2,
	\label{eq:simam-stats}
\end{equation}
and each activation is re-weighted as
\begin{equation}
	E_i^{-1}=\frac{(x_i-\mu)^2}{4(v+\lambda)}+\frac12,
	\qquad \widetilde x_i=x_i\operatorname{sigmoid}(E_i^{-1}),
	\label{eq:simam-gate}
\end{equation}
where $x_i$ is the activation at position $i$ within a channel, $M$ is the number of positions in that channel, $\mu$ and $v$ are the mean and variance of the channel, $\lambda$ is a regularisation constant, $E_i^{-1}$ is the reciprocal of the minimal energy of position $i$, and $\widetilde x_i$ is the re-weighted activation. The reciprocal energy measures distinctiveness, and the sigmoid bounds the resulting multiplicative weight while preserving the order of the scores. The printed shared-variance definition in the paper uses $M$, whereas the code in its Figure~3 uses $M-1$; the implementation follows the latter.

In this model, $M=THW$ within each channel, rather than the $HW$ of the original image formulation~\cite{yang2021}, and $\lambda=10^{-4}$. SimAM acts on each convolutional stream before concatenation. It adds no learned parameters but still performs reductions and element-wise computations. Whole-clip statistics do not guarantee that apex frames receive the highest weights, since unusual activation values may also reflect noise.

\subsection{Self-Attention and the Transformer Encoder}
\label{sec:self-attention}

In self-attention, each input token is projected into a query, a key and a value. Scaled dot-product attention~\cite{vaswani2017}, in the form also presented by Dosovitskiy et al.~\cite{dosovitskiy2021}, computes
\begin{equation}
	\operatorname{Attention}(Q,K,V)
	=\operatorname{softmax}\left(\frac{QK^{\mathsf T}}{\sqrt{d_k}}\right)V,
	\label{eq:attention}
\end{equation}
where $Q$, $K$ and $V$ are the matrices of queries, keys and values, with one row per token, $d_k$ is the dimension of the key vectors, and the softmax is applied to each row. The queries and keys determine the contribution of each value to each output. The scaling by $\sqrt{d_k}$ moderates the magnitude of the dot products as the key dimension grows. Multiple attention heads learn projections into different subspaces, whose outputs are concatenated and linearly combined~\cite{vaswani2017}.

Without position information, self-attention is permutation-equivariant: permuting the input tokens permutes their outputs correspondingly, so the temporal order is not identified. Fixed sinusoidal position vectors~\cite{vaswani2017} provide this information in the present model. Although the sinusoidal functions can be evaluated at other positions, generalisation to unseen sequence lengths is not guaranteed, and the implementation uses a finite precomputed buffer.

The temporal encoder alternates attention and position-wise feed-forward sublayers with residual connections~\cite{he2016}. It uses pre-norm placement, applying layer normalisation before each sublayer~\cite{xiong2020}. Its output is averaged over time before classification. This mean summarises temporal features that are already contextualised, and therefore differs from averaging the input tokens without an encoder.

\subsection{Replacement Operations for Disabled Components}
\label{sec:fallbacks}

\autoref{tab:ch2-component-controls} defines the operation that replaces each component when it is disabled.

\begin{table}[htbp]
    \centering
    \small
    \begin{tabular}{@{}p{0.22\linewidth} p{0.35\linewidth} p{0.35\linewidth}@{}}
        \hline
        \textbf{Component} & \textbf{Enabled} & \textbf{Disabled} \\
        \hline
        Spatial stem & Three learned convolutional streams & Per-frame $4\times4$ average pooling and a learned $48\rightarrow96$ linear projection \\
        SimAM & Per-stream activation re-weighting & Unmodified stream outputs; the model enables SimAM only together with its CNN stem \\
        Temporal encoder & Position encoding, transformer layers, then temporal mean & Direct mean of the input sequence \\
        \hline
    \end{tabular}
    \caption[Enabled components and their controls]{Operations compared by the model switches. The classification head is shared in structure; compatible feature dimensions do not imply equal parameter counts.}
    \label{tab:ch2-component-controls}
\end{table}

Pooling and direct temporal averaging have no learned weights. The projection used without the CNN stem does: including its bias, it contains $48\times96+96=4{,}704$ parameters. The replacements preserve the interfaces required by the experiment while changing the operations and the capacity. They do not isolate architecture independently of parameter count or computational cost.


\section{Training with Limited and Imbalanced Data}
\label{sec:scarcity-skew}

Micro-expression datasets are small and their emotion classes are unevenly represented. This section introduces the loss functions, sampling strategy, augmentation and optimisation methods used to train under these conditions.

\subsection{Cross-Entropy under Class Imbalance}
\label{sec:skew-cross-entropy}

The cross-entropy loss of a single example is
\begin{equation}
	L_{\mathrm{CE}}=-\log p_t,
	\label{eq:cross-entropy}
\end{equation}
where $p_t$ is the probability predicted for the true class $t$. Under natural sampling, a frequent class contributes more terms to the training objective, which biases learning towards the majority class~\cite{he2009}. A low overall loss therefore does not imply that rare classes are recognised well.

\subsection{Focal Loss}
\label{sec:focal-loss}

Focal loss, introduced by Lin et al.~\cite{lin2017} and applied to micro-expression recognition by Zhao et al.~\cite{zhaoS2021}, is defined as
\begin{equation}
	\operatorname{FL}(p_t)=-\alpha_t(1-p_t)^\gamma\log p_t,
	\label{eq:focal-loss}
\end{equation}
where $p_t$ is the predicted probability of the true class, $\gamma\geq0$ is the focusing parameter and $\alpha_t$ is an optional weighting factor for the true class. For $\gamma>0$, the modulating factor $(1-p_t)^\gamma$ reduces the contribution of confidently classified examples relative to difficult ones. This factor depends on prediction difficulty, not on class frequency; any frequency correction must come from $\alpha_t$ or from the sampling strategy. This thesis uses $\gamma=2$, the value reported by Zhao et al.~\cite{zhaoS2021}.

\subsection{Class-Balanced Sampling}
\label{sec:balanced-sampling}
\label{sec:combining-both}

Oversampling the minority classes is a common data-level response to class imbalance~\cite{he2009,buda2018}. Each training example is assigned the sampling weight
\begin{equation}
	w_i=\frac{1}{n_{c_i}},
	\label{eq:sampler-weight}
\end{equation}
where $w_i$ is the weight of training example $i$, $c_i$ is its class, and $n_{c_i}$ is the number of training examples of that class. Drawing with replacement according to these weights makes each class equally likely in expectation. Balanced sampling changes how often each class is seen, not the number of distinct examples, and it is applied to training only; evaluation uses the natural class distribution of the held-out data.

Class weights in the loss ($\alpha_t$ in \eqref{eq:focal-loss}) and balanced sampling correct the same imbalance at different points. Applying both corrects it twice. In this thesis the frequency correction is therefore assigned to the sampler alone, and the loss carries no class weights.

\subsection{Label Smoothing}
\label{sec:label-smoothing}

Label smoothing~\cite{szegedy2016} replaces the one-hot target with a mixture of the one-hot target and a uniform distribution. For $C$ classes, the smoothed target is
\begin{equation}
	q_c=(1-\epsilon)\mathbf{1}[c=t]+\frac{\epsilon}{C},
	\label{eq:smoothed-target}
\end{equation}
and the corresponding cross-entropy term is
\begin{equation}
	L_{\mathrm{smooth}}
	=-(1-\epsilon)\log p_t
	-\frac{\epsilon}{C}\sum_{c=1}^{C}\log p_c,
	\label{eq:label-smoothing}
\end{equation}
where $q_c$ is the target probability of class $c$, $\epsilon$ is the smoothing strength, $t$ is the true class, $\mathbf{1}[c=t]$ equals one when $c=t$ and zero otherwise, and $p_c$ is the predicted probability of class $c$. Label smoothing discourages over-confident predictions~\cite{szegedy2016}. In this thesis, the focal factor $(1-p_t)^\gamma$ of \eqref{eq:focal-loss} multiplies the smoothed term \eqref{eq:label-smoothing}, with $p_t$ taken from the unsmoothed prediction.

\subsection{Data Augmentation}
\label{sec:augmentation}

Data augmentation enlarges the effective training set through label-preserving transformations, and Xia et al.~\cite{xia2020a} motivate it for limited and imbalanced micro-expression data. The augmentation used here is a random horizontal flip, applied with probability one half during training only. Because mirroring a clip reverses the direction of horizontal motion, the horizontal flow channel $u$ is negated whenever the clip is flipped.

\subsection{Optimisation}
\label{sec:optimisation}

AdamW~\cite{loshchilov2019}, a variant of Adam~\cite{kingma2015}, updates the parameters with weight decay decoupled from the gradient-based update. The learning rate is increased linearly during an initial warmup phase and then reduced by cosine annealing~\cite{loshchilov2017}, a schedule also used for micro-expression recognition by Zhang et al.~\cite{zhang2022}. Mixed-precision training~\cite{micikevicius2018} performs most computations in 16-bit floating point to reduce memory use, with gradient scaling to prevent small gradient values from underflowing.


\section{Evaluation Methodology}
\label{sec:evaluation}

This section defines the metrics and the cross-validation protocol used to evaluate the models.

\subsection{Confusion Matrix, Precision, Recall and F1}
\label{sec:confusion-matrix}

In a confusion matrix, entry $(i,j)$ counts the examples whose true class is $i$ and whose predicted class is $j$, so the diagonal contains the correct predictions. The class-wise metrics follow OFF-ApexNet~\cite{liong2019a} and the MEGC 2019 protocol~\cite{see2019}:
\begin{equation}
	P_c=\frac{TP_c}{TP_c+FP_c},\qquad
	R_c=\frac{TP_c}{TP_c+FN_c},\qquad
	F1_c=\frac{2TP_c}{2TP_c+FP_c+FN_c},
	\label{eq:prf}
\end{equation}
where $P_c$, $R_c$ and $F1_c$ are the precision, recall and F1 score of class $c$; $TP_c$ (true positives) counts correct predictions of $c$; $FP_c$ (false positives) counts incorrect predictions of $c$; and $FN_c$ (false negatives) counts examples of $c$ predicted as another class. Precision penalises false alarms and recall penalises misses. F1 is their harmonic mean, so a high value requires both to be high. When a denominator is zero, the score is set to zero.

\subsection{Macro and Micro Averaging}
\label{sec:macro-micro}

Macro F1 gives each class equal weight, using the mean of per-class F1 scores specified by MEGC 2019~\cite{see2019}:
\begin{equation}
	F1_{\mathrm{macro}}=\frac1C\sum_{c=1}^{C}F1_c,
	\label{eq:macro-f1}
\end{equation}
where $C$ is the number of classes and $F1_c$ is defined in \eqref{eq:prf}. Micro F1 instead sums the true positives, false positives and false negatives across classes before computing the score~\cite{sokolova2009}. In single-label multiclass classification, micro F1 equals accuracy and can therefore be dominated by a frequent class. Macro F1 is the more informative measure under class imbalance.

\subsection{Pooled and Per-Fold Averaging}
\label{sec:pooled-vs-fold}

Under cross-validation, macro F1 can be computed in two ways: by pooling the confusion matrices of all folds before computing the class F1 scores, or by computing macro F1 in each fold and averaging the fold scores. Because F1 is non-linear, the two need not agree. If a fold contains no examples of a class, that class contributes zero to the fold score even when every example is classified correctly, so averaging fold scores mixes classifier behaviour with fold composition. Pooling avoids this. The Unweighted F1 (UF1) of MEGC 2019 is defined through the pooled counts~\cite{see2019}, and pooled macro F1 is the primary metric in this thesis.

\subsection{Cross-Validation and Subject Independence}
\label{sec:cross-validation}

Cross-validation repeatedly fits a model to a training partition and evaluates it on held-out examples. Leave-one-subject-out (LOSO) cross-validation holds out all clips of one subject in each fold and trains on the remaining subjects. It is the native protocol of CASME~II~\cite{yan2014} and the protocol mandated by MEGC 2019~\cite{see2019}. Because no subject appears in both partitions, the model cannot exploit familiarity with a subject's identity or recording conditions.

Two further principles govern how such results can be interpreted. First, if the held-out data are also used for model selection, such as choosing the best training epoch, the resulting estimate is optimistically biased; an inner validation split drawn from the training subjects avoids this~\cite{varma2006}. Second, testing whether two classifiers evaluated on the same examples differ significantly requires a paired test such as McNemar's test~\cite{mcnemar1947}, which needs the prediction for every individual example.


\section{Summary}
\label{sec:ch2-conclusion}

This chapter introduced the concepts on which the thesis builds. It defined micro-expressions and the properties that make them difficult to recognise: short duration and low intensity, and the relationship between action units and emotion labels. It described the CASME~II recordings and their preparation, Eulerian video magnification, and the optical flow and optical strain representations that form the model input. It then introduced the neural network concepts and building blocks used by the model, the methods used to train it on limited and imbalanced data, and the metrics and leave-one-subject-out protocol used to evaluate it.

\autoref{ch:literature} reviews the published work on each of these components. For each stage of the pipeline, it examines the evidence for its reported contribution to micro-expression recognition and identifies the gap that this thesis addresses.
.
% This file is the single source for Chapter 2; the older chapter2.tex and the
% chapter2_short_half*.tex files are superseded and must not be included.
%
% Equations are numbered and labelled eq:*; each is followed by a definition of
% its symbols. The mathematics of every equation is unchanged from the version
% reproduced in contents/formula_guide_complete.tex.
%
% Figure 2.1 (fig2_1_casme2_sequence.png) is Figure 2 of Yan et al. (2014),
% PLoS ONE 9(1):e86041, reproduced under the CC BY licence with attribution.
%
% Requires in the preamble: amsmath, graphicx.
% The pipeline overview (formerly the last section) now lives in Chapter 4.
% =============================================================================

\chapter{Background}
\label{ch:background}

This chapter presents the concepts on which the remainder of the thesis builds. \autoref{sec:phenomenon} defines micro-expressions and their characteristics. \autoref{sec:recording-to-input} describes the recordings of the CASME~II dataset and their preparation as model input. \autoref{sec:motion-magnification} introduces Eulerian video magnification, and \autoref{sec:video-to-motion} the optical flow and optical strain representations derived from the magnified video. \autoref{sec:nn-fundamentals} and \autoref{sec:building-blocks} cover the neural network concepts and building blocks used by the model. \autoref{sec:scarcity-skew} addresses training with limited and imbalanced data, and \autoref{sec:evaluation} presents the evaluation methodology.


\section{Micro-Expressions}
\label{sec:phenomenon}

Facial expressions are commonly divided into macro-expressions and micro-expressions. Macro-expressions are the ordinary expressions of everyday interaction; they typically last between 0.5 and 4 seconds and involve a large area of the face~\cite{ben2021,xie2022,adegun2020}. Micro-expressions are brief, involuntary and local facial movements that tend to occur when a person attempts to conceal or suppress an emotion, particularly in high-stakes situations~\cite{yan2014,ben2021,xie2022,li2022survey}. Ekman and Friesen described such movements as a form of nonverbal leakage~\cite{ekman1969}, and they have since been studied empirically in concealed and falsified emotional expressions~\cite{porter2008}. Their involuntary character distinguishes them from expressions deliberately posed for a camera~\cite{qu2016}. Comprehensive reviews of automatic micro-expression analysis are provided by Oh et al.~\cite{oh2018}, Ben et al.~\cite{ben2021}, Li et al.~\cite{li2022survey} and Xie et al.~\cite{xie2022}.

\autoref{fig:casme2-sequence} shows an example from CASME~II: a lowering of the inner brow (AU~4, together with AU~9) labelled as disgust, whose apex occurs approximately 110~ms after onset~\cite{yan2014}. The change between frames is barely visible at full-face scale, which illustrates why the phenomenon is difficult to perceive. An observed movement is nevertheless neither proof of deception nor an unambiguous indicator of emotion; the recognition task depends on the annotations of the dataset.

\begin{figure}[htbp]
	\centering
	\includegraphics[width=\textwidth]{fig2_1_casme2_sequence.png}
	\caption[Frame sequence of a CASME II micro-expression]{Frame sequence of a micro-expression from CASME~II, with time in milliseconds from onset. The action units are AU~4 and AU~9, labelled as disgust; the apex occurs at approximately 110~ms. The insets magnify the right inner brow (AU~4). Reproduced from Yan et al.~\cite{yan2014} (Figure~2) under the Creative Commons Attribution (CC~BY) licence.}
	\label{fig:casme2-sequence}
\end{figure}

\subsection{Duration}
\label{sec:duration}

Definitions of micro-expression duration vary. Yan et al.~\cite{yan2014} describe half a second as the generally accepted upper limit, whereas other sources report typical durations of 1/25 to 1/3 of a second~\cite{li2022survey} or 0.065 to 0.5 seconds~\cite{ben2021}. The limit is a working definition rather than a duration filter applied by this pipeline. A brief event yields few observations at ordinary video frame rates, so capture frequency and temporal sampling determine the evidence available for recognition.

\subsection{Intensity}
\label{sec:low-intensity}

Suppressed facial movements may be difficult to perceive even when sufficient frames are recorded~\cite{yan2014}. Because micro-expressions are local~\cite{ben2021}, a small contraction may affect only the region around the brow or mouth while most of the image remains unchanged. Duration limits the number of observations, whereas low intensity limits how clearly the movement can be separated from noise within them. Both properties motivate motion magnification and explicit motion features.

\subsection{Temporal Phases: Onset, Apex and Offset}
\label{sec:onset-apex-offset}

The \textbf{onset} marks the beginning of the movement, the \textbf{apex} its point of greatest annotated intensity, and the \textbf{offset} its end (\autoref{fig:casme2-sequence}). Apex-frame methods use the peak frame as a compact description~\cite{liY2018}. This pipeline instead uses the annotated onset and offset to bound a sequence. The apex annotation is neither read by the model nor used to select frames, so temporal features must be learned from the sampled sequence rather than supplied through a selected peak.

\subsection{Action Units and the Prediction Target}
\label{sec:action-units}

The Facial Action Coding System (FACS)~\cite{ekman1978} describes facial movements through Action Units (AUs), each corresponding to an individual component of muscle movement~\cite{li2022survey}, rather than through emotion categories. CASME~II assigns emotion labels using coded movements together with other annotation evidence~\cite{yan2014}, and the mapping is not one-to-one. This thesis predicts the resulting emotion category only. AU metadata is retained in the coding table but is neither an input nor a training target, and the system does not perform AU recognition.


\section{Data Acquisition and Preprocessing}
\label{sec:recording-to-input}

Input preparation determines which recorded information reaches every subsequent component. In particular, image resizing must be distinguished from face alignment and from temporal interpolation.

\subsection{High-Speed Recording}
\label{sec:high-speed-capture}

The number of frame intervals spanned by a facial movement is approximately
\begin{equation}
	n \approx D f_s,
	\label{eq:frame-intervals}
\end{equation}
where $n$ is the number of frame intervals, $D$ is the duration of the movement in seconds and $f_s$ is the recording rate in frames per second (fps). For an event lasting $D = 0.3$~s, this gives roughly eight to nine intervals at 25--30~fps, against sixty at 200~fps. CASME~II is recorded at 200~fps~\cite{yan2014}. The pipeline takes this rate from its configuration rather than reading it from each clip. After temporal subsampling, the physical spacing between selected frames need not match the nominal rate.

\subsection{Source Frames}
\label{sec:corpus-supplies}

CASME~II provides both the original recorded frames and a registered, cropped release. The study manifest identifies folders of the original $640\times480$ frames rather than the normalised release. The preprocessing performs no face detection, landmark localisation, registration, alignment or cropping of its own. The input therefore includes the full recorded field of view, and neither facial position nor scale is geometrically standardised. The registration described by Yan et al.~\cite{yan2014} must not be attributed to this pipeline.

\subsection{Frame Preparation}
\label{sec:frame-preparation}

Frames are loaded in greyscale and resized by bilinear interpolation to $224\times224$ pixels. Resizing the whole rectangular frame to a square changes its aspect ratio; the result is not a cropped face region. The loader selects the frames between the annotated onset and offset and then samples a fixed-length sequence. No apex localisation is performed.

\subsection{Clip Selection}
\label{sec:clip-selection}

Preprocessing selects CASME~II entries that are coded as micro-expressions, marked as having frames, and have an inclusive onset-to-offset length greater than two frames. The loader must still find readable frames, since the metadata flag alone does not guarantee successful extraction. The recognition dataset additionally requires a matching saved tensor and a label in the selected three-class mapping; the residual \texttt{Others} category is excluded. Preprocessing eligibility and final classification membership are therefore separate filters.


\section{Eulerian Video Magnification}
\label{sec:motion-magnification}

\subsection{The Eulerian Approach}
\label{sec:eulerian-idea}

Eulerian video magnification (EVM), introduced by Wu et al.~\cite{wu2012}, examines intensity changes at fixed image locations rather than explicitly tracking moving points. At each pixel, it isolates temporal variation within a selected frequency band, scales that variation and adds it back to the signal. Bai et al.~\cite{bai2021} describe the amplitude-based mechanism through the small-motion approximation
\begin{equation}
	I(x,t)=f(x+\delta(t))
	\approx f(x)+\delta(t)\frac{\partial f}{\partial x},
	\label{eq:evm-small-motion}
\end{equation}
where $I(x,t)$ is the image intensity at position $x$ and time $t$, $f$ is the intensity profile of the image at the initial time, and $\delta(t)$ is the displacement of that profile at time $t$. Adding $\alpha$ times the motion-related variation gives approximately
\begin{equation}
	\widetilde I(x,t)\approx
	f(x)+(1+\alpha)\delta(t)\frac{\partial f}{\partial x},
	\label{eq:evm-magnified}
\end{equation}
where $\widetilde I(x,t)$ is the magnified intensity and $\alpha$ is the amplification factor. For sufficiently small displacements, intensity amplification therefore resembles a displacement enlarged by the factor $(1+\alpha)$. The approximation does not imply that arbitrary intensity changes represent facial movement, nor that larger amplification always improves recognition.

\subsection{Spatial Decomposition with the Laplacian Pyramid}
\label{sec:laplacian-pyramid}

A Laplacian pyramid~\cite{burt1983} separates an image into spatial scales. At each level, the image is blurred and downsampled, and an upsampled reconstruction is subtracted from the current image to retain the detail at that scale. Repeating this procedure produces a set of Laplacian detail bands and a final coarse Gaussian residual. Reconstruction starts from the residual, upsamples it and adds the detail bands successively. The implementation uses the separable blur kernel $[1,4,6,4,1]/16$ and factor-two downsampling.

\subsection{Temporal Filtering}
\label{sec:temporal-filtering}

Each pixel in a spatial band supplies a time series. A band-pass filter retains the variation within the selected temporal-frequency interval, and the retained signal is amplified before reconstruction. Bai et al.~\cite{bai2021} describe a Butterworth filter for amplitude-based magnification. The present implementation instead sets to zero the Fourier coefficients outside the selected band and then applies the inverse transform.

\subsection{The Amplification Trade-Off}
\label{sec:amplification-tradeoff}

The filter distinguishes frequencies, not causes. Head movement, illumination variation and noise within the pass-band may be amplified together with the facial contraction of interest. A larger $\alpha$ can also violate the small-displacement approximation of \eqref{eq:evm-small-motion} and produce artefacts. Magnification therefore changes the signal presented to flow estimation without guaranteeing a cleaner or more discriminative representation.

\subsection{Implementation Details}
\label{sec:evm-departures}

\autoref{tab:ch2-evm-departures} summarises the implementation choices that affect the interpretation of the magnified input.

\begin{table}[htbp]
    \centering
    \small
    \begin{tabular}{@{}p{0.26\linewidth} p{0.68\linewidth}@{}}
        \hline
        \textbf{Property} & \textbf{Implemented behaviour and consequence} \\
        \hline
        Temporal filter & Hard Fourier mask rather than the Butterworth filter described by Bai et al.~\cite{bai2021}; the pass-band has no smooth roll-off. \\
        Spatial amplification & The same $\alpha$ is used for every Laplacian band, without spatial-frequency damping. The Gaussian residual is neither filtered nor amplified. \\
        Colour treatment & Processing operates on greyscale intensity, with no separate chromatic attenuation. \\
        Pipeline ordering & Each clip is sampled to 33 frames before magnification, but the filter still assumes the nominal 200 fps. Its physical pass-band is therefore clip-dependent. \\
        \hline
    \end{tabular}
    \caption[EVM implementation departures]{Implementation choices that distinguish the magnifier and its position in the pipeline from the reference formulation. These choices are shared by all magnification-enabled configurations.}
    \label{tab:ch2-evm-departures}
\end{table}

The Laplacian pyramid itself is retained. After reconstruction, values are clipped to $[0,255]$ and converted to 8-bit greyscale before motion extraction. Because subsampling changes the frame spacing, the nominal band does not correspond to the same physical-frequency interval for every clip.


\section{Motion Representation: Optical Flow and Optical Strain}
\label{sec:video-to-motion}

\subsection{Motivation for Motion-Based Input}
\label{sec:why-motion}

Optical flow makes displacement explicit instead of requiring the network to infer it from appearance alone. This can reduce reliance on raw image intensity but does not guarantee invariance to identity or illumination: facial geometry, registration errors and brightness changes still affect the estimated field. The resulting representation therefore remains conditioned on the capture and preprocessing procedures.

\subsection{Optical Flow and Brightness Constancy}
\label{sec:optical-flow}

Optical flow estimates the apparent displacement of image content between frames~\cite{horn1981}. It rests on the brightness-constancy assumption that a moving point retains its intensity, as described for micro-expression recognition in OFF-ApexNet~\cite{liong2019a}:
\begin{equation}
	I(x,y,t)=I(x+u\Delta t,y+v\Delta t,t+\Delta t),
	\label{eq:brightness-constancy}
\end{equation}
where $I(x,y,t)$ is the intensity at pixel $(x,y)$ and time $t$, $u$ and $v$ are the horizontal and vertical velocity components, and $\Delta t$ is the time interval between frames. A first-order Taylor expansion gives the optical flow constraint equation~\cite{horn1981}
\begin{equation}
	I_xu+I_yv+I_t=0,
	\label{eq:flow-constraint}
\end{equation}
where $I_x$, $I_y$ and $I_t$ are the partial derivatives of the intensity with respect to $x$, $y$ and $t$. This is a single equation in two unknowns. Under the \textbf{aperture problem}, a local edge constrains the motion normal to the edge, along the intensity gradient, but does not determine the motion parallel to it. Estimators therefore introduce neighbourhood assumptions or regularisation. In the discrete pipeline, $u$ and $v$ denote displacement in pixels between selected frames rather than physical velocity per second. Illumination changes and large inter-frame displacements can violate the underlying approximation.

\subsection{Optical Strain}
\label{sec:optical-strain}

Optical strain describes how neighbouring displacements differ. For the displacement field $\mathbf{u}=[u,v]^{\mathsf T}$, the small-strain approximation used for optical strain~\cite{shreve2011} is
\begin{equation}
	\varepsilon=\tfrac12\left(\nabla\mathbf{u}
	+(\nabla\mathbf{u})^{\mathsf T}\right),
	\qquad
	\varepsilon_{xx}=\frac{\partial u}{\partial x},\quad
	\varepsilon_{yy}=\frac{\partial v}{\partial y},
	\label{eq:strain-tensor}
\end{equation}
\begin{equation}
	\varepsilon_{xy}=\varepsilon_{yx}
	=\tfrac12\left(\frac{\partial u}{\partial y}
	+\frac{\partial v}{\partial x}\right),
	\label{eq:strain-shear}
\end{equation}
where $\varepsilon$ is the $2\times2$ strain tensor, $\nabla\mathbf{u}$ is the spatial gradient (Jacobian) of the displacement field, $\varepsilon_{xx}$ and $\varepsilon_{yy}$ are the normal strain components in the horizontal and vertical directions, and $\varepsilon_{xy}=\varepsilon_{yx}$ is the shear strain component. The normal components describe local stretching or compression, and the off-diagonal terms describe shear. A spatially uniform translation has zero gradient and therefore zero strain. This property does not imply immunity to arbitrary head rotation, projection effects or flow-estimation errors, and differentiation may amplify noise.

The implemented scalar strain magnitude is
\begin{equation}
	\varepsilon_{\mathrm{mag}}
	=\sqrt{\varepsilon_{xx}^2+\varepsilon_{yy}^2
		+\varepsilon_{xy}^2},
	\label{eq:strain-magnitude}
\end{equation}
where $\varepsilon_{\mathrm{mag}}$ is the strain magnitude at a pixel and the remaining symbols are as defined in \eqref{eq:strain-tensor} and \eqref{eq:strain-shear}. This form counts shear once. The Frobenius norm of the symmetric tensor counts both off-diagonal entries, giving $2\varepsilon_{xy}^2$ inside the root. The two forms differ in the relative weighting of shear, not merely by a removable global scale. The scalar discards the direction of deformation but accompanies the original displacement components as a third channel.

\subsection{Construction of the Three-Channel Input Tensor}
\label{sec:three-channel-tensor}

Uniform index sampling selects 33 existing frames from the onset-to-offset sequence. No Temporal Interpolation Model (TIM)~\cite{pfister2011} or other frame synthesis is performed, and repeated indices can occur for short clips. After optional magnification, Farneb\"ack optical flow~\cite{farneback2003} is computed for the 32 adjacent frame pairs, following the estimator choice reported by Zhao et al.~\cite{zhaoS2021}. The OpenCV settings are pyramid scale 0.5, three levels, window size 15, three iterations, polynomial neighbourhood 5, polynomial smoothing 1.2 and flags 0.

The spatial derivatives required for strain are computed with NumPy's \texttt{gradient} function, using unit pixel spacing, central differences in the interior and first-order one-sided differences at the boundaries. Temporal derivatives are not used, and no dedicated strain post-filter is applied. Stacking $[u,v,\varepsilon_{\mathrm{mag}}]$ gives a tensor of shape $(3,32,224,224)$ over the resized full frame. The network receives this tensor rather than the original intensity sequence.

Two normalisations are applied. During extraction, each channel is mapped independently to $[0,1]$ using its extrema across the clip. During loading, each channel is standardised over its $T\times H\times W$ values, after any training augmentation. Small constants in the denominators protect nearly constant channels, which therefore need not have unit variance. Both normalisations operate per sample and use no statistics pooled across training and held-out subjects.


\section{Neural Network Fundamentals}
\label{sec:nn-fundamentals}

A neural network composes parameterised operations to map the motion tensor to emotion-class scores~\cite{lecun2015,goodfellow2016}. This section introduces only the concepts required to interpret the implemented blocks and their controls.

\subsection{Units, Layers and the Forward Pass}
\label{sec:units-layers}

A unit computes
\begin{equation}
	a=\phi(w^{\mathsf T}x+b),
	\label{eq:unit}
\end{equation}
where $x$ is the input vector, $w$ is the learned weight vector, $b$ is the learned bias, $\phi$ is an activation function and $a$ is the resulting activation~\cite{goodfellow2016}. A layer applies several such operations in parallel, and composing layers gives a network. Depth refers to the number of layers and width to their feature dimension. A forward pass evaluates these operations with the current parameters and produces class scores, or logits, from which class probabilities are obtained.

\subsection{Non-Linear Activation Functions}
\label{sec:non-linearity}

Without non-linear operations, a stack of affine layers reduces to a single affine map. Non-linearity allows the composition to represent more complex relationships. The convolutional streams use the rectified linear unit (ReLU)~\cite{nair2010}, and the transformer feed-forward layers use the Gaussian error linear unit (GELU)~\cite{hendrycks2016}:
\begin{equation}
	\operatorname{ReLU}(x)=\max(0,x),
	\qquad
	\operatorname{GELU}(x)=x\,\Phi(x),
	\label{eq:activations}
\end{equation}
where $x$ is the input to the activation and $\Phi$ is the cumulative distribution function of the standard normal distribution. Both are fixed functions, although they act on features learned by the surrounding layers.

\subsection{Gradient-Based Learning}
\label{sec:gradient-descent}

A loss function compares predictions with labels. Backpropagation~\cite{rumelhart1986} applies the chain rule through the network to compute the derivatives of the loss with respect to its parameters. An optimiser uses these gradients to update the parameters, with the learning rate controlling the size of each update. Mini-batches provide successive updates, and an epoch completes the configured number of batch draws. When sampling is performed with replacement, an epoch need not visit every distinct clip.

During evaluation, the parameters are held fixed and training-only stochastic behaviour such as dropout is disabled. A low training loss alone does not establish useful learning, because the relevant objective is prediction on unseen subjects. A validation set used to select an epoch therefore has a different role from a test set reserved for final assessment.

\subsection{Convolution}
\label{sec:convolution}

A convolution applies the same learned kernel at multiple spatial positions~\cite{lecun1998}. This weight sharing allows local patterns to be detected without learning a separate detector for every image location. Each output channel is a feature map, and each kernel combines all input channels over its local window. For a two-dimensional layer with one bias per output channel, the number of parameters is
\begin{equation}
	N_{\mathrm{params}}=C_{\mathrm{out}}
	(C_{\mathrm{in}}k_Hk_W+1),
	\label{eq:conv-params}
\end{equation}
where $C_{\mathrm{in}}$ and $C_{\mathrm{out}}$ are the numbers of input and output channels, $k_H$ and $k_W$ are the kernel height and width, and the added one accounts for the bias of each output channel. The parameter count is independent of the image dimensions, but computation and activation memory increase with the number of positions processed. The stride sets the spacing between kernel applications, and padding controls boundary treatment and output size. Stacking layers enlarges the receptive field, that is, the input region that can influence a later activation.

\subsection{Pooling}
\label{sec:pooling}

Pooling replaces a neighbourhood with a fixed summary such as its maximum or mean~\cite{lecun1998,goodfellow2016}. A $2\times2$ max-pooling operation with stride two halves the spatial dimensions and retains the strongest response in each window. Adaptive average pooling instead produces a specified output grid. Neither operation learns weights. Both discard detail while reducing the computation of subsequent layers; any later learned projection is a separate operation with its own parameters.

\subsection{Model Capacity and Small Datasets}
\label{sec:capacity}

Capacity describes the range of functions a model can represent~\cite{goodfellow2016}. The parameter count is one indicator, but architecture, constraints and training procedure also matter. A model can overfit subject-specific appearance, noise or recording conditions, achieving a low training loss without generalising. Having more parameters than examples does not by itself prove overfitting, just as few parameters do not guarantee good predictions.

The small size of the CASME~II training pool motivates controlled architectural comparisons and regularisation. Their effectiveness must be assessed on held-out subjects, with checkpoint selection kept separate from final testing where possible. Because the analytical flow and strain transformations already expose motion structure, the learned network does not start from raw pixels.


\section{Network Building Blocks}
\label{sec:building-blocks}

The model combines a spatial stem, optional activation re-weighting, an optional temporal encoder and a common classification head. The stem consists of three streams with separate learned weights, one each for horizontal flow, vertical flow and strain. Disabling a component replaces a specified operation while preserving compatible input and output shapes.

\subsection{Three-Dimensional Convolution and Kernel Shape}
\label{sec:conv3d}

A three-dimensional convolution~\cite{ji2013,tran2015} operates on tensors of shape $(B,C,T,H,W)$, denoting batch size, channels, time, height and width. A kernel of shape $(k_T,k_H,k_W)$ mixes information across time only when $k_T>1$. The implemented streams use $(1,3,3)$ kernels with $(0,1,1)$ padding, followed by $(1,2,2)$ max-pooling. The convolutions therefore preserve the temporal length and learn spatial filters, while pooling halves the height and width.

The use of \texttt{Conv3d} does not establish that temporal convolution was tested. The complete stream can nevertheless depend on multiple frames, because batch normalisation during training and SimAM, whenever enabled, use statistics that span time. This distinction separates temporal convolution from other forms of whole-clip dependence.

\subsection{Normalisation and Regularisation}
\label{sec:normalisation}

Batch normalisation~\cite{ioffe2015} rescales each channel using a mean and variance, followed by a learned scale and shift:
\begin{equation}
	\widehat x=\frac{x-\mu_c}{\sqrt{\sigma_c^2+\epsilon}},
	\qquad y=\gamma_c\widehat x+\beta_c,
	\label{eq:batchnorm}
\end{equation}
where $x$ is an activation in channel $c$, $\mu_c$ and $\sigma_c^2$ are the mean and variance of that channel, $\epsilon$ is a small constant for numerical stability, $\widehat x$ is the normalised activation, $\gamma_c$ and $\beta_c$ are the learned scale and shift of the channel, and $y$ is the output. \texttt{BatchNorm3d} computes the training statistics over the $(B,T,H,W)$ dimensions of each channel, and evaluation uses running estimates accumulated during training. These activation statistics differ from the per-sample input standardisation performed by the data loader.

Dropout~\cite{srivastava2014} randomly deactivates units during training to reduce overfitting. \texttt{Dropout3d} removes entire feature channels for a sample~\cite{tompson2015} rather than masking individual spatial values. The retained channels are rescaled, and evaluation disables the random mask. Layer normalisation~\cite{ba2016} instead standardises the feature dimension of each token independently of other samples. It appears inside the transformer and in the classification head, where it is followed by dropout and a linear projection to class logits.

\subsection{Parameter-Free Attention: SimAM}
\label{sec:simam}

SimAM assigns larger weights to activations that differ from their channel context~\cite{yang2021}. Following the reference code of the original paper, the statistics of each context are
\begin{equation}
	\mu=\frac1M\sum_i x_i,\qquad
	v=\frac1{M-1}\sum_i(x_i-\mu)^2,
	\label{eq:simam-stats}
\end{equation}
and each activation is re-weighted as
\begin{equation}
	E_i^{-1}=\frac{(x_i-\mu)^2}{4(v+\lambda)}+\frac12,
	\qquad \widetilde x_i=x_i\operatorname{sigmoid}(E_i^{-1}),
	\label{eq:simam-gate}
\end{equation}
where $x_i$ is the activation at position $i$ within a channel, $M$ is the number of positions in that channel, $\mu$ and $v$ are the mean and variance of the channel, $\lambda$ is a regularisation constant, $E_i^{-1}$ is the reciprocal of the minimal energy of position $i$, and $\widetilde x_i$ is the re-weighted activation. The reciprocal energy measures distinctiveness, and the sigmoid bounds the resulting multiplicative weight while preserving the order of the scores. The printed shared-variance definition in the paper uses $M$, whereas the code in its Figure~3 uses $M-1$; the implementation follows the latter.

In this model, $M=THW$ within each channel, rather than the $HW$ of the original image formulation~\cite{yang2021}, and $\lambda=10^{-4}$. SimAM acts on each convolutional stream before concatenation. It adds no learned parameters but still performs reductions and element-wise computations. Whole-clip statistics do not guarantee that apex frames receive the highest weights, since unusual activation values may also reflect noise.

\subsection{Self-Attention and the Transformer Encoder}
\label{sec:self-attention}

In self-attention, each input token is projected into a query, a key and a value. Scaled dot-product attention~\cite{vaswani2017}, in the form also presented by Dosovitskiy et al.~\cite{dosovitskiy2021}, computes
\begin{equation}
	\operatorname{Attention}(Q,K,V)
	=\operatorname{softmax}\left(\frac{QK^{\mathsf T}}{\sqrt{d_k}}\right)V,
	\label{eq:attention}
\end{equation}
where $Q$, $K$ and $V$ are the matrices of queries, keys and values, with one row per token, $d_k$ is the dimension of the key vectors, and the softmax is applied to each row. The queries and keys determine the contribution of each value to each output. The scaling by $\sqrt{d_k}$ moderates the magnitude of the dot products as the key dimension grows. Multiple attention heads learn projections into different subspaces, whose outputs are concatenated and linearly combined~\cite{vaswani2017}.

Without position information, self-attention is permutation-equivariant: permuting the input tokens permutes their outputs correspondingly, so the temporal order is not identified. Fixed sinusoidal position vectors~\cite{vaswani2017} provide this information in the present model. Although the sinusoidal functions can be evaluated at other positions, generalisation to unseen sequence lengths is not guaranteed, and the implementation uses a finite precomputed buffer.

The temporal encoder alternates attention and position-wise feed-forward sublayers with residual connections~\cite{he2016}. It uses pre-norm placement, applying layer normalisation before each sublayer~\cite{xiong2020}. Its output is averaged over time before classification. This mean summarises temporal features that are already contextualised, and therefore differs from averaging the input tokens without an encoder.

\subsection{Replacement Operations for Disabled Components}
\label{sec:fallbacks}

\autoref{tab:ch2-component-controls} defines the operation that replaces each component when it is disabled.

\begin{table}[htbp]
    \centering
    \small
    \begin{tabular}{@{}p{0.22\linewidth} p{0.35\linewidth} p{0.35\linewidth}@{}}
        \hline
        \textbf{Component} & \textbf{Enabled} & \textbf{Disabled} \\
        \hline
        Spatial stem & Three learned convolutional streams & Per-frame $4\times4$ average pooling and a learned $48\rightarrow96$ linear projection \\
        SimAM & Per-stream activation re-weighting & Unmodified stream outputs; the model enables SimAM only together with its CNN stem \\
        Temporal encoder & Position encoding, transformer layers, then temporal mean & Direct mean of the input sequence \\
        \hline
    \end{tabular}
    \caption[Enabled components and their controls]{Operations compared by the model switches. The classification head is shared in structure; compatible feature dimensions do not imply equal parameter counts.}
    \label{tab:ch2-component-controls}
\end{table}

Pooling and direct temporal averaging have no learned weights. The projection used without the CNN stem does: including its bias, it contains $48\times96+96=4{,}704$ parameters. The replacements preserve the interfaces required by the experiment while changing the operations and the capacity. They do not isolate architecture independently of parameter count or computational cost.


\section{Training with Limited and Imbalanced Data}
\label{sec:scarcity-skew}

Micro-expression datasets are small and their emotion classes are unevenly represented. This section introduces the loss functions, sampling strategy, augmentation and optimisation methods used to train under these conditions.

\subsection{Cross-Entropy under Class Imbalance}
\label{sec:skew-cross-entropy}

The cross-entropy loss of a single example is
\begin{equation}
	L_{\mathrm{CE}}=-\log p_t,
	\label{eq:cross-entropy}
\end{equation}
where $p_t$ is the probability predicted for the true class $t$. Under natural sampling, a frequent class contributes more terms to the training objective, which biases learning towards the majority class~\cite{he2009}. A low overall loss therefore does not imply that rare classes are recognised well.

\subsection{Focal Loss}
\label{sec:focal-loss}

Focal loss, introduced by Lin et al.~\cite{lin2017} and applied to micro-expression recognition by Zhao et al.~\cite{zhaoS2021}, is defined as
\begin{equation}
	\operatorname{FL}(p_t)=-\alpha_t(1-p_t)^\gamma\log p_t,
	\label{eq:focal-loss}
\end{equation}
where $p_t$ is the predicted probability of the true class, $\gamma\geq0$ is the focusing parameter and $\alpha_t$ is an optional weighting factor for the true class. For $\gamma>0$, the modulating factor $(1-p_t)^\gamma$ reduces the contribution of confidently classified examples relative to difficult ones. This factor depends on prediction difficulty, not on class frequency; any frequency correction must come from $\alpha_t$ or from the sampling strategy. This thesis uses $\gamma=2$, the value reported by Zhao et al.~\cite{zhaoS2021}.

\subsection{Class-Balanced Sampling}
\label{sec:balanced-sampling}
\label{sec:combining-both}

Oversampling the minority classes is a common data-level response to class imbalance~\cite{he2009,buda2018}. Each training example is assigned the sampling weight
\begin{equation}
	w_i=\frac{1}{n_{c_i}},
	\label{eq:sampler-weight}
\end{equation}
where $w_i$ is the weight of training example $i$, $c_i$ is its class, and $n_{c_i}$ is the number of training examples of that class. Drawing with replacement according to these weights makes each class equally likely in expectation. Balanced sampling changes how often each class is seen, not the number of distinct examples, and it is applied to training only; evaluation uses the natural class distribution of the held-out data.

Class weights in the loss ($\alpha_t$ in \eqref{eq:focal-loss}) and balanced sampling correct the same imbalance at different points. Applying both corrects it twice. In this thesis the frequency correction is therefore assigned to the sampler alone, and the loss carries no class weights.

\subsection{Label Smoothing}
\label{sec:label-smoothing}

Label smoothing~\cite{szegedy2016} replaces the one-hot target with a mixture of the one-hot target and a uniform distribution. For $C$ classes, the smoothed target is
\begin{equation}
	q_c=(1-\epsilon)\mathbf{1}[c=t]+\frac{\epsilon}{C},
	\label{eq:smoothed-target}
\end{equation}
and the corresponding cross-entropy term is
\begin{equation}
	L_{\mathrm{smooth}}
	=-(1-\epsilon)\log p_t
	-\frac{\epsilon}{C}\sum_{c=1}^{C}\log p_c,
	\label{eq:label-smoothing}
\end{equation}
where $q_c$ is the target probability of class $c$, $\epsilon$ is the smoothing strength, $t$ is the true class, $\mathbf{1}[c=t]$ equals one when $c=t$ and zero otherwise, and $p_c$ is the predicted probability of class $c$. Label smoothing discourages over-confident predictions~\cite{szegedy2016}. In this thesis, the focal factor $(1-p_t)^\gamma$ of \eqref{eq:focal-loss} multiplies the smoothed term \eqref{eq:label-smoothing}, with $p_t$ taken from the unsmoothed prediction.

\subsection{Data Augmentation}
\label{sec:augmentation}

Data augmentation enlarges the effective training set through label-preserving transformations, and Xia et al.~\cite{xia2020a} motivate it for limited and imbalanced micro-expression data. The augmentation used here is a random horizontal flip, applied with probability one half during training only. Because mirroring a clip reverses the direction of horizontal motion, the horizontal flow channel $u$ is negated whenever the clip is flipped.

\subsection{Optimisation}
\label{sec:optimisation}

AdamW~\cite{loshchilov2019}, a variant of Adam~\cite{kingma2015}, updates the parameters with weight decay decoupled from the gradient-based update. The learning rate is increased linearly during an initial warmup phase and then reduced by cosine annealing~\cite{loshchilov2017}, a schedule also used for micro-expression recognition by Zhang et al.~\cite{zhang2022}. Mixed-precision training~\cite{micikevicius2018} performs most computations in 16-bit floating point to reduce memory use, with gradient scaling to prevent small gradient values from underflowing.


\section{Evaluation Methodology}
\label{sec:evaluation}

This section defines the metrics and the cross-validation protocol used to evaluate the models.

\subsection{Confusion Matrix, Precision, Recall and F1}
\label{sec:confusion-matrix}

In a confusion matrix, entry $(i,j)$ counts the examples whose true class is $i$ and whose predicted class is $j$, so the diagonal contains the correct predictions. The class-wise metrics follow OFF-ApexNet~\cite{liong2019a} and the MEGC 2019 protocol~\cite{see2019}:
\begin{equation}
	P_c=\frac{TP_c}{TP_c+FP_c},\qquad
	R_c=\frac{TP_c}{TP_c+FN_c},\qquad
	F1_c=\frac{2TP_c}{2TP_c+FP_c+FN_c},
	\label{eq:prf}
\end{equation}
where $P_c$, $R_c$ and $F1_c$ are the precision, recall and F1 score of class $c$; $TP_c$ (true positives) counts correct predictions of $c$; $FP_c$ (false positives) counts incorrect predictions of $c$; and $FN_c$ (false negatives) counts examples of $c$ predicted as another class. Precision penalises false alarms and recall penalises misses. F1 is their harmonic mean, so a high value requires both to be high. When a denominator is zero, the score is set to zero.

\subsection{Macro and Micro Averaging}
\label{sec:macro-micro}

Macro F1 gives each class equal weight, using the mean of per-class F1 scores specified by MEGC 2019~\cite{see2019}:
\begin{equation}
	F1_{\mathrm{macro}}=\frac1C\sum_{c=1}^{C}F1_c,
	\label{eq:macro-f1}
\end{equation}
where $C$ is the number of classes and $F1_c$ is defined in \eqref{eq:prf}. Micro F1 instead sums the true positives, false positives and false negatives across classes before computing the score~\cite{sokolova2009}. In single-label multiclass classification, micro F1 equals accuracy and can therefore be dominated by a frequent class. Macro F1 is the more informative measure under class imbalance.

\subsection{Pooled and Per-Fold Averaging}
\label{sec:pooled-vs-fold}

Under cross-validation, macro F1 can be computed in two ways: by pooling the confusion matrices of all folds before computing the class F1 scores, or by computing macro F1 in each fold and averaging the fold scores. Because F1 is non-linear, the two need not agree. If a fold contains no examples of a class, that class contributes zero to the fold score even when every example is classified correctly, so averaging fold scores mixes classifier behaviour with fold composition. Pooling avoids this. The Unweighted F1 (UF1) of MEGC 2019 is defined through the pooled counts~\cite{see2019}, and pooled macro F1 is the primary metric in this thesis.

\subsection{Cross-Validation and Subject Independence}
\label{sec:cross-validation}

Cross-validation repeatedly fits a model to a training partition and evaluates it on held-out examples. Leave-one-subject-out (LOSO) cross-validation holds out all clips of one subject in each fold and trains on the remaining subjects. It is the native protocol of CASME~II~\cite{yan2014} and the protocol mandated by MEGC 2019~\cite{see2019}. Because no subject appears in both partitions, the model cannot exploit familiarity with a subject's identity or recording conditions.

Two further principles govern how such results can be interpreted. First, if the held-out data are also used for model selection, such as choosing the best training epoch, the resulting estimate is optimistically biased; an inner validation split drawn from the training subjects avoids this~\cite{varma2006}. Second, testing whether two classifiers evaluated on the same examples differ significantly requires a paired test such as McNemar's test~\cite{mcnemar1947}, which needs the prediction for every individual example.


\section{Summary}
\label{sec:ch2-conclusion}

This chapter introduced the concepts on which the thesis builds. It defined micro-expressions and the properties that make them difficult to recognise: short duration and low intensity, and the relationship between action units and emotion labels. It described the CASME~II recordings and their preparation, Eulerian video magnification, and the optical flow and optical strain representations that form the model input. It then introduced the neural network concepts and building blocks used by the model, the methods used to train it on limited and imbalanced data, and the metrics and leave-one-subject-out protocol used to evaluate it.

\autoref{ch:literature} reviews the published work on each of these components. For each stage of the pipeline, it examines the evidence for its reported contribution to micro-expression recognition and identifies the gap that this thesis addresses.
.
% This file is the single source for Chapter 2; the older chapter2.tex and the
% chapter2_short_half*.tex files are superseded and must not be included.
%
% Equations are numbered and labelled eq:*; each is followed by a definition of
% its symbols. The mathematics of every equation is unchanged from the version
% reproduced in contents/formula_guide_complete.tex.
%
% Figure 2.1 (fig2_1_casme2_sequence.png) is Figure 2 of Yan et al. (2014),
% PLoS ONE 9(1):e86041, reproduced under the CC BY licence with attribution.
%
% Requires in the preamble: amsmath, graphicx.
% The pipeline overview (formerly the last section) now lives in Chapter 4.
% =============================================================================

\chapter{Background}
\label{ch:background}

This chapter presents the concepts on which the remainder of the thesis builds. \autoref{sec:phenomenon} defines micro-expressions and their characteristics. \autoref{sec:recording-to-input} describes the recordings of the CASME~II dataset and their preparation as model input. \autoref{sec:motion-magnification} introduces Eulerian video magnification, and \autoref{sec:video-to-motion} the optical flow and optical strain representations derived from the magnified video. \autoref{sec:nn-fundamentals} and \autoref{sec:building-blocks} cover the neural network concepts and building blocks used by the model. \autoref{sec:scarcity-skew} addresses training with limited and imbalanced data, and \autoref{sec:evaluation} presents the evaluation methodology.


\section{Micro-Expressions}
\label{sec:phenomenon}

Facial expressions are commonly divided into macro-expressions and micro-expressions. Macro-expressions are the ordinary expressions of everyday interaction; they typically last between 0.5 and 4 seconds and involve a large area of the face~\cite{ben2021,xie2022,adegun2020}. Micro-expressions are brief, involuntary and local facial movements that tend to occur when a person attempts to conceal or suppress an emotion, particularly in high-stakes situations~\cite{yan2014,ben2021,xie2022,li2022survey}. Ekman and Friesen described such movements as a form of nonverbal leakage~\cite{ekman1969}, and they have since been studied empirically in concealed and falsified emotional expressions~\cite{porter2008}. Their involuntary character distinguishes them from expressions deliberately posed for a camera~\cite{qu2016}. Comprehensive reviews of automatic micro-expression analysis are provided by Oh et al.~\cite{oh2018}, Ben et al.~\cite{ben2021}, Li et al.~\cite{li2022survey} and Xie et al.~\cite{xie2022}.

\autoref{fig:casme2-sequence} shows an example from CASME~II: a lowering of the inner brow (AU~4, together with AU~9) labelled as disgust, whose apex occurs approximately 110~ms after onset~\cite{yan2014}. The change between frames is barely visible at full-face scale, which illustrates why the phenomenon is difficult to perceive. An observed movement is nevertheless neither proof of deception nor an unambiguous indicator of emotion; the recognition task depends on the annotations of the dataset.

\begin{figure}[htbp]
	\centering
	\includegraphics[width=\textwidth]{fig2_1_casme2_sequence.png}
	\caption[Frame sequence of a CASME II micro-expression]{Frame sequence of a micro-expression from CASME~II, with time in milliseconds from onset. The action units are AU~4 and AU~9, labelled as disgust; the apex occurs at approximately 110~ms. The insets magnify the right inner brow (AU~4). Reproduced from Yan et al.~\cite{yan2014} (Figure~2) under the Creative Commons Attribution (CC~BY) licence.}
	\label{fig:casme2-sequence}
\end{figure}

\subsection{Duration}
\label{sec:duration}

Definitions of micro-expression duration vary. Yan et al.~\cite{yan2014} describe half a second as the generally accepted upper limit, whereas other sources report typical durations of 1/25 to 1/3 of a second~\cite{li2022survey} or 0.065 to 0.5 seconds~\cite{ben2021}. The limit is a working definition rather than a duration filter applied by this pipeline. A brief event yields few observations at ordinary video frame rates, so capture frequency and temporal sampling determine the evidence available for recognition.

\subsection{Intensity}
\label{sec:low-intensity}

Suppressed facial movements may be difficult to perceive even when sufficient frames are recorded~\cite{yan2014}. Because micro-expressions are local~\cite{ben2021}, a small contraction may affect only the region around the brow or mouth while most of the image remains unchanged. Duration limits the number of observations, whereas low intensity limits how clearly the movement can be separated from noise within them. Both properties motivate motion magnification and explicit motion features.

\subsection{Temporal Phases: Onset, Apex and Offset}
\label{sec:onset-apex-offset}

The \textbf{onset} marks the beginning of the movement, the \textbf{apex} its point of greatest annotated intensity, and the \textbf{offset} its end (\autoref{fig:casme2-sequence}). Apex-frame methods use the peak frame as a compact description~\cite{liY2018}. This pipeline instead uses the annotated onset and offset to bound a sequence. The apex annotation is neither read by the model nor used to select frames, so temporal features must be learned from the sampled sequence rather than supplied through a selected peak.

\subsection{Action Units and the Prediction Target}
\label{sec:action-units}

The Facial Action Coding System (FACS)~\cite{ekman1978} describes facial movements through Action Units (AUs), each corresponding to an individual component of muscle movement~\cite{li2022survey}, rather than through emotion categories. CASME~II assigns emotion labels using coded movements together with other annotation evidence~\cite{yan2014}, and the mapping is not one-to-one. This thesis predicts the resulting emotion category only. AU metadata is retained in the coding table but is neither an input nor a training target, and the system does not perform AU recognition.


\section{Data Acquisition and Preprocessing}
\label{sec:recording-to-input}

Input preparation determines which recorded information reaches every subsequent component. In particular, image resizing must be distinguished from face alignment and from temporal interpolation.

\subsection{High-Speed Recording}
\label{sec:high-speed-capture}

The number of frame intervals spanned by a facial movement is approximately
\begin{equation}
	n \approx D f_s,
	\label{eq:frame-intervals}
\end{equation}
where $n$ is the number of frame intervals, $D$ is the duration of the movement in seconds and $f_s$ is the recording rate in frames per second (fps). For an event lasting $D = 0.3$~s, this gives roughly eight to nine intervals at 25--30~fps, against sixty at 200~fps. CASME~II is recorded at 200~fps~\cite{yan2014}. The pipeline takes this rate from its configuration rather than reading it from each clip. After temporal subsampling, the physical spacing between selected frames need not match the nominal rate.

\subsection{Source Frames}
\label{sec:corpus-supplies}

CASME~II provides both the original recorded frames and a registered, cropped release. The study manifest identifies folders of the original $640\times480$ frames rather than the normalised release. The preprocessing performs no face detection, landmark localisation, registration, alignment or cropping of its own. The input therefore includes the full recorded field of view, and neither facial position nor scale is geometrically standardised. The registration described by Yan et al.~\cite{yan2014} must not be attributed to this pipeline.

\subsection{Frame Preparation}
\label{sec:frame-preparation}

Frames are loaded in greyscale and resized by bilinear interpolation to $224\times224$ pixels. Resizing the whole rectangular frame to a square changes its aspect ratio; the result is not a cropped face region. The loader selects the frames between the annotated onset and offset and then samples a fixed-length sequence. No apex localisation is performed.

\subsection{Clip Selection}
\label{sec:clip-selection}

Preprocessing selects CASME~II entries that are coded as micro-expressions, marked as having frames, and have an inclusive onset-to-offset length greater than two frames. The loader must still find readable frames, since the metadata flag alone does not guarantee successful extraction. The recognition dataset additionally requires a matching saved tensor and a label in the selected three-class mapping; the residual \texttt{Others} category is excluded. Preprocessing eligibility and final classification membership are therefore separate filters.


\section{Eulerian Video Magnification}
\label{sec:motion-magnification}

\subsection{The Eulerian Approach}
\label{sec:eulerian-idea}

Eulerian video magnification (EVM), introduced by Wu et al.~\cite{wu2012}, examines intensity changes at fixed image locations rather than explicitly tracking moving points. At each pixel, it isolates temporal variation within a selected frequency band, scales that variation and adds it back to the signal. Bai et al.~\cite{bai2021} describe the amplitude-based mechanism through the small-motion approximation
\begin{equation}
	I(x,t)=f(x+\delta(t))
	\approx f(x)+\delta(t)\frac{\partial f}{\partial x},
	\label{eq:evm-small-motion}
\end{equation}
where $I(x,t)$ is the image intensity at position $x$ and time $t$, $f$ is the intensity profile of the image at the initial time, and $\delta(t)$ is the displacement of that profile at time $t$. Adding $\alpha$ times the motion-related variation gives approximately
\begin{equation}
	\widetilde I(x,t)\approx
	f(x)+(1+\alpha)\delta(t)\frac{\partial f}{\partial x},
	\label{eq:evm-magnified}
\end{equation}
where $\widetilde I(x,t)$ is the magnified intensity and $\alpha$ is the amplification factor. For sufficiently small displacements, intensity amplification therefore resembles a displacement enlarged by the factor $(1+\alpha)$. The approximation does not imply that arbitrary intensity changes represent facial movement, nor that larger amplification always improves recognition.

\subsection{Spatial Decomposition with the Laplacian Pyramid}
\label{sec:laplacian-pyramid}

A Laplacian pyramid~\cite{burt1983} separates an image into spatial scales. At each level, the image is blurred and downsampled, and an upsampled reconstruction is subtracted from the current image to retain the detail at that scale. Repeating this procedure produces a set of Laplacian detail bands and a final coarse Gaussian residual. Reconstruction starts from the residual, upsamples it and adds the detail bands successively. The implementation uses the separable blur kernel $[1,4,6,4,1]/16$ and factor-two downsampling.

\subsection{Temporal Filtering}
\label{sec:temporal-filtering}

Each pixel in a spatial band supplies a time series. A band-pass filter retains the variation within the selected temporal-frequency interval, and the retained signal is amplified before reconstruction. Bai et al.~\cite{bai2021} describe a Butterworth filter for amplitude-based magnification. The present implementation instead sets to zero the Fourier coefficients outside the selected band and then applies the inverse transform.

\subsection{The Amplification Trade-Off}
\label{sec:amplification-tradeoff}

The filter distinguishes frequencies, not causes. Head movement, illumination variation and noise within the pass-band may be amplified together with the facial contraction of interest. A larger $\alpha$ can also violate the small-displacement approximation of \eqref{eq:evm-small-motion} and produce artefacts. Magnification therefore changes the signal presented to flow estimation without guaranteeing a cleaner or more discriminative representation.

\subsection{Implementation Details}
\label{sec:evm-departures}

\autoref{tab:ch2-evm-departures} summarises the implementation choices that affect the interpretation of the magnified input.

\begin{table}[htbp]
    \centering
    \small
    \begin{tabular}{@{}p{0.26\linewidth} p{0.68\linewidth}@{}}
        \hline
        \textbf{Property} & \textbf{Implemented behaviour and consequence} \\
        \hline
        Temporal filter & Hard Fourier mask rather than the Butterworth filter described by Bai et al.~\cite{bai2021}; the pass-band has no smooth roll-off. \\
        Spatial amplification & The same $\alpha$ is used for every Laplacian band, without spatial-frequency damping. The Gaussian residual is neither filtered nor amplified. \\
        Colour treatment & Processing operates on greyscale intensity, with no separate chromatic attenuation. \\
        Pipeline ordering & Each clip is sampled to 33 frames before magnification, but the filter still assumes the nominal 200 fps. Its physical pass-band is therefore clip-dependent. \\
        \hline
    \end{tabular}
    \caption[EVM implementation departures]{Implementation choices that distinguish the magnifier and its position in the pipeline from the reference formulation. These choices are shared by all magnification-enabled configurations.}
    \label{tab:ch2-evm-departures}
\end{table}

The Laplacian pyramid itself is retained. After reconstruction, values are clipped to $[0,255]$ and converted to 8-bit greyscale before motion extraction. Because subsampling changes the frame spacing, the nominal band does not correspond to the same physical-frequency interval for every clip.


\section{Motion Representation: Optical Flow and Optical Strain}
\label{sec:video-to-motion}

\subsection{Motivation for Motion-Based Input}
\label{sec:why-motion}

Optical flow makes displacement explicit instead of requiring the network to infer it from appearance alone. This can reduce reliance on raw image intensity but does not guarantee invariance to identity or illumination: facial geometry, registration errors and brightness changes still affect the estimated field. The resulting representation therefore remains conditioned on the capture and preprocessing procedures.

\subsection{Optical Flow and Brightness Constancy}
\label{sec:optical-flow}

Optical flow estimates the apparent displacement of image content between frames~\cite{horn1981}. It rests on the brightness-constancy assumption that a moving point retains its intensity, as described for micro-expression recognition in OFF-ApexNet~\cite{liong2019a}:
\begin{equation}
	I(x,y,t)=I(x+u\Delta t,y+v\Delta t,t+\Delta t),
	\label{eq:brightness-constancy}
\end{equation}
where $I(x,y,t)$ is the intensity at pixel $(x,y)$ and time $t$, $u$ and $v$ are the horizontal and vertical velocity components, and $\Delta t$ is the time interval between frames. A first-order Taylor expansion gives the optical flow constraint equation~\cite{horn1981}
\begin{equation}
	I_xu+I_yv+I_t=0,
	\label{eq:flow-constraint}
\end{equation}
where $I_x$, $I_y$ and $I_t$ are the partial derivatives of the intensity with respect to $x$, $y$ and $t$. This is a single equation in two unknowns. Under the \textbf{aperture problem}, a local edge constrains the motion normal to the edge, along the intensity gradient, but does not determine the motion parallel to it. Estimators therefore introduce neighbourhood assumptions or regularisation. In the discrete pipeline, $u$ and $v$ denote displacement in pixels between selected frames rather than physical velocity per second. Illumination changes and large inter-frame displacements can violate the underlying approximation.

\subsection{Optical Strain}
\label{sec:optical-strain}

Optical strain describes how neighbouring displacements differ. For the displacement field $\mathbf{u}=[u,v]^{\mathsf T}$, the small-strain approximation used for optical strain~\cite{shreve2011} is
\begin{equation}
	\varepsilon=\tfrac12\left(\nabla\mathbf{u}
	+(\nabla\mathbf{u})^{\mathsf T}\right),
	\qquad
	\varepsilon_{xx}=\frac{\partial u}{\partial x},\quad
	\varepsilon_{yy}=\frac{\partial v}{\partial y},
	\label{eq:strain-tensor}
\end{equation}
\begin{equation}
	\varepsilon_{xy}=\varepsilon_{yx}
	=\tfrac12\left(\frac{\partial u}{\partial y}
	+\frac{\partial v}{\partial x}\right),
	\label{eq:strain-shear}
\end{equation}
where $\varepsilon$ is the $2\times2$ strain tensor, $\nabla\mathbf{u}$ is the spatial gradient (Jacobian) of the displacement field, $\varepsilon_{xx}$ and $\varepsilon_{yy}$ are the normal strain components in the horizontal and vertical directions, and $\varepsilon_{xy}=\varepsilon_{yx}$ is the shear strain component. The normal components describe local stretching or compression, and the off-diagonal terms describe shear. A spatially uniform translation has zero gradient and therefore zero strain. This property does not imply immunity to arbitrary head rotation, projection effects or flow-estimation errors, and differentiation may amplify noise.

The implemented scalar strain magnitude is
\begin{equation}
	\varepsilon_{\mathrm{mag}}
	=\sqrt{\varepsilon_{xx}^2+\varepsilon_{yy}^2
		+\varepsilon_{xy}^2},
	\label{eq:strain-magnitude}
\end{equation}
where $\varepsilon_{\mathrm{mag}}$ is the strain magnitude at a pixel and the remaining symbols are as defined in \eqref{eq:strain-tensor} and \eqref{eq:strain-shear}. This form counts shear once. The Frobenius norm of the symmetric tensor counts both off-diagonal entries, giving $2\varepsilon_{xy}^2$ inside the root. The two forms differ in the relative weighting of shear, not merely by a removable global scale. The scalar discards the direction of deformation but accompanies the original displacement components as a third channel.

\subsection{Construction of the Three-Channel Input Tensor}
\label{sec:three-channel-tensor}

Uniform index sampling selects 33 existing frames from the onset-to-offset sequence. No Temporal Interpolation Model (TIM)~\cite{pfister2011} or other frame synthesis is performed, and repeated indices can occur for short clips. After optional magnification, Farneb\"ack optical flow~\cite{farneback2003} is computed for the 32 adjacent frame pairs, following the estimator choice reported by Zhao et al.~\cite{zhaoS2021}. The OpenCV settings are pyramid scale 0.5, three levels, window size 15, three iterations, polynomial neighbourhood 5, polynomial smoothing 1.2 and flags 0.

The spatial derivatives required for strain are computed with NumPy's \texttt{gradient} function, using unit pixel spacing, central differences in the interior and first-order one-sided differences at the boundaries. Temporal derivatives are not used, and no dedicated strain post-filter is applied. Stacking $[u,v,\varepsilon_{\mathrm{mag}}]$ gives a tensor of shape $(3,32,224,224)$ over the resized full frame. The network receives this tensor rather than the original intensity sequence.

Two normalisations are applied. During extraction, each channel is mapped independently to $[0,1]$ using its extrema across the clip. During loading, each channel is standardised over its $T\times H\times W$ values, after any training augmentation. Small constants in the denominators protect nearly constant channels, which therefore need not have unit variance. Both normalisations operate per sample and use no statistics pooled across training and held-out subjects.


\section{Neural Network Fundamentals}
\label{sec:nn-fundamentals}

A neural network composes parameterised operations to map the motion tensor to emotion-class scores~\cite{lecun2015,goodfellow2016}. This section introduces only the concepts required to interpret the implemented blocks and their controls.

\subsection{Units, Layers and the Forward Pass}
\label{sec:units-layers}

A unit computes
\begin{equation}
	a=\phi(w^{\mathsf T}x+b),
	\label{eq:unit}
\end{equation}
where $x$ is the input vector, $w$ is the learned weight vector, $b$ is the learned bias, $\phi$ is an activation function and $a$ is the resulting activation~\cite{goodfellow2016}. A layer applies several such operations in parallel, and composing layers gives a network. Depth refers to the number of layers and width to their feature dimension. A forward pass evaluates these operations with the current parameters and produces class scores, or logits, from which class probabilities are obtained.

\subsection{Non-Linear Activation Functions}
\label{sec:non-linearity}

Without non-linear operations, a stack of affine layers reduces to a single affine map. Non-linearity allows the composition to represent more complex relationships. The convolutional streams use the rectified linear unit (ReLU)~\cite{nair2010}, and the transformer feed-forward layers use the Gaussian error linear unit (GELU)~\cite{hendrycks2016}:
\begin{equation}
	\operatorname{ReLU}(x)=\max(0,x),
	\qquad
	\operatorname{GELU}(x)=x\,\Phi(x),
	\label{eq:activations}
\end{equation}
where $x$ is the input to the activation and $\Phi$ is the cumulative distribution function of the standard normal distribution. Both are fixed functions, although they act on features learned by the surrounding layers.

\subsection{Gradient-Based Learning}
\label{sec:gradient-descent}

A loss function compares predictions with labels. Backpropagation~\cite{rumelhart1986} applies the chain rule through the network to compute the derivatives of the loss with respect to its parameters. An optimiser uses these gradients to update the parameters, with the learning rate controlling the size of each update. Mini-batches provide successive updates, and an epoch completes the configured number of batch draws. When sampling is performed with replacement, an epoch need not visit every distinct clip.

During evaluation, the parameters are held fixed and training-only stochastic behaviour such as dropout is disabled. A low training loss alone does not establish useful learning, because the relevant objective is prediction on unseen subjects. A validation set used to select an epoch therefore has a different role from a test set reserved for final assessment.

\subsection{Convolution}
\label{sec:convolution}

A convolution applies the same learned kernel at multiple spatial positions~\cite{lecun1998}. This weight sharing allows local patterns to be detected without learning a separate detector for every image location. Each output channel is a feature map, and each kernel combines all input channels over its local window. For a two-dimensional layer with one bias per output channel, the number of parameters is
\begin{equation}
	N_{\mathrm{params}}=C_{\mathrm{out}}
	(C_{\mathrm{in}}k_Hk_W+1),
	\label{eq:conv-params}
\end{equation}
where $C_{\mathrm{in}}$ and $C_{\mathrm{out}}$ are the numbers of input and output channels, $k_H$ and $k_W$ are the kernel height and width, and the added one accounts for the bias of each output channel. The parameter count is independent of the image dimensions, but computation and activation memory increase with the number of positions processed. The stride sets the spacing between kernel applications, and padding controls boundary treatment and output size. Stacking layers enlarges the receptive field, that is, the input region that can influence a later activation.

\subsection{Pooling}
\label{sec:pooling}

Pooling replaces a neighbourhood with a fixed summary such as its maximum or mean~\cite{lecun1998,goodfellow2016}. A $2\times2$ max-pooling operation with stride two halves the spatial dimensions and retains the strongest response in each window. Adaptive average pooling instead produces a specified output grid. Neither operation learns weights. Both discard detail while reducing the computation of subsequent layers; any later learned projection is a separate operation with its own parameters.

\subsection{Model Capacity and Small Datasets}
\label{sec:capacity}

Capacity describes the range of functions a model can represent~\cite{goodfellow2016}. The parameter count is one indicator, but architecture, constraints and training procedure also matter. A model can overfit subject-specific appearance, noise or recording conditions, achieving a low training loss without generalising. Having more parameters than examples does not by itself prove overfitting, just as few parameters do not guarantee good predictions.

The small size of the CASME~II training pool motivates controlled architectural comparisons and regularisation. Their effectiveness must be assessed on held-out subjects, with checkpoint selection kept separate from final testing where possible. Because the analytical flow and strain transformations already expose motion structure, the learned network does not start from raw pixels.


\section{Network Building Blocks}
\label{sec:building-blocks}

The model combines a spatial stem, optional activation re-weighting, an optional temporal encoder and a common classification head. The stem consists of three streams with separate learned weights, one each for horizontal flow, vertical flow and strain. Disabling a component replaces a specified operation while preserving compatible input and output shapes.

\subsection{Three-Dimensional Convolution and Kernel Shape}
\label{sec:conv3d}

A three-dimensional convolution~\cite{ji2013,tran2015} operates on tensors of shape $(B,C,T,H,W)$, denoting batch size, channels, time, height and width. A kernel of shape $(k_T,k_H,k_W)$ mixes information across time only when $k_T>1$. The implemented streams use $(1,3,3)$ kernels with $(0,1,1)$ padding, followed by $(1,2,2)$ max-pooling. The convolutions therefore preserve the temporal length and learn spatial filters, while pooling halves the height and width.

The use of \texttt{Conv3d} does not establish that temporal convolution was tested. The complete stream can nevertheless depend on multiple frames, because batch normalisation during training and SimAM, whenever enabled, use statistics that span time. This distinction separates temporal convolution from other forms of whole-clip dependence.

\subsection{Normalisation and Regularisation}
\label{sec:normalisation}

Batch normalisation~\cite{ioffe2015} rescales each channel using a mean and variance, followed by a learned scale and shift:
\begin{equation}
	\widehat x=\frac{x-\mu_c}{\sqrt{\sigma_c^2+\epsilon}},
	\qquad y=\gamma_c\widehat x+\beta_c,
	\label{eq:batchnorm}
\end{equation}
where $x$ is an activation in channel $c$, $\mu_c$ and $\sigma_c^2$ are the mean and variance of that channel, $\epsilon$ is a small constant for numerical stability, $\widehat x$ is the normalised activation, $\gamma_c$ and $\beta_c$ are the learned scale and shift of the channel, and $y$ is the output. \texttt{BatchNorm3d} computes the training statistics over the $(B,T,H,W)$ dimensions of each channel, and evaluation uses running estimates accumulated during training. These activation statistics differ from the per-sample input standardisation performed by the data loader.

Dropout~\cite{srivastava2014} randomly deactivates units during training to reduce overfitting. \texttt{Dropout3d} removes entire feature channels for a sample~\cite{tompson2015} rather than masking individual spatial values. The retained channels are rescaled, and evaluation disables the random mask. Layer normalisation~\cite{ba2016} instead standardises the feature dimension of each token independently of other samples. It appears inside the transformer and in the classification head, where it is followed by dropout and a linear projection to class logits.

\subsection{Parameter-Free Attention: SimAM}
\label{sec:simam}

SimAM assigns larger weights to activations that differ from their channel context~\cite{yang2021}. Following the reference code of the original paper, the statistics of each context are
\begin{equation}
	\mu=\frac1M\sum_i x_i,\qquad
	v=\frac1{M-1}\sum_i(x_i-\mu)^2,
	\label{eq:simam-stats}
\end{equation}
and each activation is re-weighted as
\begin{equation}
	E_i^{-1}=\frac{(x_i-\mu)^2}{4(v+\lambda)}+\frac12,
	\qquad \widetilde x_i=x_i\operatorname{sigmoid}(E_i^{-1}),
	\label{eq:simam-gate}
\end{equation}
where $x_i$ is the activation at position $i$ within a channel, $M$ is the number of positions in that channel, $\mu$ and $v$ are the mean and variance of the channel, $\lambda$ is a regularisation constant, $E_i^{-1}$ is the reciprocal of the minimal energy of position $i$, and $\widetilde x_i$ is the re-weighted activation. The reciprocal energy measures distinctiveness, and the sigmoid bounds the resulting multiplicative weight while preserving the order of the scores. The printed shared-variance definition in the paper uses $M$, whereas the code in its Figure~3 uses $M-1$; the implementation follows the latter.

In this model, $M=THW$ within each channel, rather than the $HW$ of the original image formulation~\cite{yang2021}, and $\lambda=10^{-4}$. SimAM acts on each convolutional stream before concatenation. It adds no learned parameters but still performs reductions and element-wise computations. Whole-clip statistics do not guarantee that apex frames receive the highest weights, since unusual activation values may also reflect noise.

\subsection{Self-Attention and the Transformer Encoder}
\label{sec:self-attention}

In self-attention, each input token is projected into a query, a key and a value. Scaled dot-product attention~\cite{vaswani2017}, in the form also presented by Dosovitskiy et al.~\cite{dosovitskiy2021}, computes
\begin{equation}
	\operatorname{Attention}(Q,K,V)
	=\operatorname{softmax}\left(\frac{QK^{\mathsf T}}{\sqrt{d_k}}\right)V,
	\label{eq:attention}
\end{equation}
where $Q$, $K$ and $V$ are the matrices of queries, keys and values, with one row per token, $d_k$ is the dimension of the key vectors, and the softmax is applied to each row. The queries and keys determine the contribution of each value to each output. The scaling by $\sqrt{d_k}$ moderates the magnitude of the dot products as the key dimension grows. Multiple attention heads learn projections into different subspaces, whose outputs are concatenated and linearly combined~\cite{vaswani2017}.

Without position information, self-attention is permutation-equivariant: permuting the input tokens permutes their outputs correspondingly, so the temporal order is not identified. Fixed sinusoidal position vectors~\cite{vaswani2017} provide this information in the present model. Although the sinusoidal functions can be evaluated at other positions, generalisation to unseen sequence lengths is not guaranteed, and the implementation uses a finite precomputed buffer.

The temporal encoder alternates attention and position-wise feed-forward sublayers with residual connections~\cite{he2016}. It uses pre-norm placement, applying layer normalisation before each sublayer~\cite{xiong2020}. Its output is averaged over time before classification. This mean summarises temporal features that are already contextualised, and therefore differs from averaging the input tokens without an encoder.

\subsection{Replacement Operations for Disabled Components}
\label{sec:fallbacks}

\autoref{tab:ch2-component-controls} defines the operation that replaces each component when it is disabled.

\begin{table}[htbp]
    \centering
    \small
    \begin{tabular}{@{}p{0.22\linewidth} p{0.35\linewidth} p{0.35\linewidth}@{}}
        \hline
        \textbf{Component} & \textbf{Enabled} & \textbf{Disabled} \\
        \hline
        Spatial stem & Three learned convolutional streams & Per-frame $4\times4$ average pooling and a learned $48\rightarrow96$ linear projection \\
        SimAM & Per-stream activation re-weighting & Unmodified stream outputs; the model enables SimAM only together with its CNN stem \\
        Temporal encoder & Position encoding, transformer layers, then temporal mean & Direct mean of the input sequence \\
        \hline
    \end{tabular}
    \caption[Enabled components and their controls]{Operations compared by the model switches. The classification head is shared in structure; compatible feature dimensions do not imply equal parameter counts.}
    \label{tab:ch2-component-controls}
\end{table}

Pooling and direct temporal averaging have no learned weights. The projection used without the CNN stem does: including its bias, it contains $48\times96+96=4{,}704$ parameters. The replacements preserve the interfaces required by the experiment while changing the operations and the capacity. They do not isolate architecture independently of parameter count or computational cost.


\section{Training with Limited and Imbalanced Data}
\label{sec:scarcity-skew}

Micro-expression datasets are small and their emotion classes are unevenly represented. This section introduces the loss functions, sampling strategy, augmentation and optimisation methods used to train under these conditions.

\subsection{Cross-Entropy under Class Imbalance}
\label{sec:skew-cross-entropy}

The cross-entropy loss of a single example is
\begin{equation}
	L_{\mathrm{CE}}=-\log p_t,
	\label{eq:cross-entropy}
\end{equation}
where $p_t$ is the probability predicted for the true class $t$. Under natural sampling, a frequent class contributes more terms to the training objective, which biases learning towards the majority class~\cite{he2009}. A low overall loss therefore does not imply that rare classes are recognised well.

\subsection{Focal Loss}
\label{sec:focal-loss}

Focal loss, introduced by Lin et al.~\cite{lin2017} and applied to micro-expression recognition by Zhao et al.~\cite{zhaoS2021}, is defined as
\begin{equation}
	\operatorname{FL}(p_t)=-\alpha_t(1-p_t)^\gamma\log p_t,
	\label{eq:focal-loss}
\end{equation}
where $p_t$ is the predicted probability of the true class, $\gamma\geq0$ is the focusing parameter and $\alpha_t$ is an optional weighting factor for the true class. For $\gamma>0$, the modulating factor $(1-p_t)^\gamma$ reduces the contribution of confidently classified examples relative to difficult ones. This factor depends on prediction difficulty, not on class frequency; any frequency correction must come from $\alpha_t$ or from the sampling strategy. This thesis uses $\gamma=2$, the value reported by Zhao et al.~\cite{zhaoS2021}.

\subsection{Class-Balanced Sampling}
\label{sec:balanced-sampling}
\label{sec:combining-both}

Oversampling the minority classes is a common data-level response to class imbalance~\cite{he2009,buda2018}. Each training example is assigned the sampling weight
\begin{equation}
	w_i=\frac{1}{n_{c_i}},
	\label{eq:sampler-weight}
\end{equation}
where $w_i$ is the weight of training example $i$, $c_i$ is its class, and $n_{c_i}$ is the number of training examples of that class. Drawing with replacement according to these weights makes each class equally likely in expectation. Balanced sampling changes how often each class is seen, not the number of distinct examples, and it is applied to training only; evaluation uses the natural class distribution of the held-out data.

Class weights in the loss ($\alpha_t$ in \eqref{eq:focal-loss}) and balanced sampling correct the same imbalance at different points. Applying both corrects it twice. In this thesis the frequency correction is therefore assigned to the sampler alone, and the loss carries no class weights.

\subsection{Label Smoothing}
\label{sec:label-smoothing}

Label smoothing~\cite{szegedy2016} replaces the one-hot target with a mixture of the one-hot target and a uniform distribution. For $C$ classes, the smoothed target is
\begin{equation}
	q_c=(1-\epsilon)\mathbf{1}[c=t]+\frac{\epsilon}{C},
	\label{eq:smoothed-target}
\end{equation}
and the corresponding cross-entropy term is
\begin{equation}
	L_{\mathrm{smooth}}
	=-(1-\epsilon)\log p_t
	-\frac{\epsilon}{C}\sum_{c=1}^{C}\log p_c,
	\label{eq:label-smoothing}
\end{equation}
where $q_c$ is the target probability of class $c$, $\epsilon$ is the smoothing strength, $t$ is the true class, $\mathbf{1}[c=t]$ equals one when $c=t$ and zero otherwise, and $p_c$ is the predicted probability of class $c$. Label smoothing discourages over-confident predictions~\cite{szegedy2016}. In this thesis, the focal factor $(1-p_t)^\gamma$ of \eqref{eq:focal-loss} multiplies the smoothed term \eqref{eq:label-smoothing}, with $p_t$ taken from the unsmoothed prediction.

\subsection{Data Augmentation}
\label{sec:augmentation}

Data augmentation enlarges the effective training set through label-preserving transformations, and Xia et al.~\cite{xia2020a} motivate it for limited and imbalanced micro-expression data. The augmentation used here is a random horizontal flip, applied with probability one half during training only. Because mirroring a clip reverses the direction of horizontal motion, the horizontal flow channel $u$ is negated whenever the clip is flipped.

\subsection{Optimisation}
\label{sec:optimisation}

AdamW~\cite{loshchilov2019}, a variant of Adam~\cite{kingma2015}, updates the parameters with weight decay decoupled from the gradient-based update. The learning rate is increased linearly during an initial warmup phase and then reduced by cosine annealing~\cite{loshchilov2017}, a schedule also used for micro-expression recognition by Zhang et al.~\cite{zhang2022}. Mixed-precision training~\cite{micikevicius2018} performs most computations in 16-bit floating point to reduce memory use, with gradient scaling to prevent small gradient values from underflowing.


\section{Evaluation Methodology}
\label{sec:evaluation}

This section defines the metrics and the cross-validation protocol used to evaluate the models.

\subsection{Confusion Matrix, Precision, Recall and F1}
\label{sec:confusion-matrix}

In a confusion matrix, entry $(i,j)$ counts the examples whose true class is $i$ and whose predicted class is $j$, so the diagonal contains the correct predictions. The class-wise metrics follow OFF-ApexNet~\cite{liong2019a} and the MEGC 2019 protocol~\cite{see2019}:
\begin{equation}
	P_c=\frac{TP_c}{TP_c+FP_c},\qquad
	R_c=\frac{TP_c}{TP_c+FN_c},\qquad
	F1_c=\frac{2TP_c}{2TP_c+FP_c+FN_c},
	\label{eq:prf}
\end{equation}
where $P_c$, $R_c$ and $F1_c$ are the precision, recall and F1 score of class $c$; $TP_c$ (true positives) counts correct predictions of $c$; $FP_c$ (false positives) counts incorrect predictions of $c$; and $FN_c$ (false negatives) counts examples of $c$ predicted as another class. Precision penalises false alarms and recall penalises misses. F1 is their harmonic mean, so a high value requires both to be high. When a denominator is zero, the score is set to zero.

\subsection{Macro and Micro Averaging}
\label{sec:macro-micro}

Macro F1 gives each class equal weight, using the mean of per-class F1 scores specified by MEGC 2019~\cite{see2019}:
\begin{equation}
	F1_{\mathrm{macro}}=\frac1C\sum_{c=1}^{C}F1_c,
	\label{eq:macro-f1}
\end{equation}
where $C$ is the number of classes and $F1_c$ is defined in \eqref{eq:prf}. Micro F1 instead sums the true positives, false positives and false negatives across classes before computing the score~\cite{sokolova2009}. In single-label multiclass classification, micro F1 equals accuracy and can therefore be dominated by a frequent class. Macro F1 is the more informative measure under class imbalance.

\subsection{Pooled and Per-Fold Averaging}
\label{sec:pooled-vs-fold}

Under cross-validation, macro F1 can be computed in two ways: by pooling the confusion matrices of all folds before computing the class F1 scores, or by computing macro F1 in each fold and averaging the fold scores. Because F1 is non-linear, the two need not agree. If a fold contains no examples of a class, that class contributes zero to the fold score even when every example is classified correctly, so averaging fold scores mixes classifier behaviour with fold composition. Pooling avoids this. The Unweighted F1 (UF1) of MEGC 2019 is defined through the pooled counts~\cite{see2019}, and pooled macro F1 is the primary metric in this thesis.

\subsection{Cross-Validation and Subject Independence}
\label{sec:cross-validation}

Cross-validation repeatedly fits a model to a training partition and evaluates it on held-out examples. Leave-one-subject-out (LOSO) cross-validation holds out all clips of one subject in each fold and trains on the remaining subjects. It is the native protocol of CASME~II~\cite{yan2014} and the protocol mandated by MEGC 2019~\cite{see2019}. Because no subject appears in both partitions, the model cannot exploit familiarity with a subject's identity or recording conditions.

Two further principles govern how such results can be interpreted. First, if the held-out data are also used for model selection, such as choosing the best training epoch, the resulting estimate is optimistically biased; an inner validation split drawn from the training subjects avoids this~\cite{varma2006}. Second, testing whether two classifiers evaluated on the same examples differ significantly requires a paired test such as McNemar's test~\cite{mcnemar1947}, which needs the prediction for every individual example.


\section{Summary}
\label{sec:ch2-conclusion}

This chapter introduced the concepts on which the thesis builds. It defined micro-expressions and the properties that make them difficult to recognise: short duration and low intensity, and the relationship between action units and emotion labels. It described the CASME~II recordings and their preparation, Eulerian video magnification, and the optical flow and optical strain representations that form the model input. It then introduced the neural network concepts and building blocks used by the model, the methods used to train it on limited and imbalanced data, and the metrics and leave-one-subject-out protocol used to evaluate it.

\autoref{ch:literature} reviews the published work on each of these components. For each stage of the pipeline, it examines the evidence for its reported contribution to micro-expression recognition and identifies the gap that this thesis addresses.
.
% This file is the single source for Chapter 2; the older chapter2.tex and the
% chapter2_short_half*.tex files are superseded and must not be included.
%
% Equations are numbered and labelled eq:*; each is followed by a definition of
% its symbols. The mathematics of every equation is unchanged from the version
% reproduced in contents/formula_guide_complete.tex.
%
% Figure 2.1 (fig2_1_casme2_sequence.png) is Figure 2 of Yan et al. (2014),
% PLoS ONE 9(1):e86041, reproduced under the CC BY licence with attribution.
%
% Requires in the preamble: amsmath, graphicx.
% The pipeline overview (formerly the last section) now lives in Chapter 4.
% =============================================================================

\chapter{Background}
\label{ch:background}

This chapter presents the concepts on which the remainder of the thesis builds. \autoref{sec:phenomenon} defines micro-expressions and their characteristics. \autoref{sec:recording-to-input} describes the recordings of the CASME~II dataset and their preparation as model input. \autoref{sec:motion-magnification} introduces Eulerian video magnification, and \autoref{sec:video-to-motion} the optical flow and optical strain representations derived from the magnified video. \autoref{sec:nn-fundamentals} and \autoref{sec:building-blocks} cover the neural network concepts and building blocks used by the model. \autoref{sec:scarcity-skew} addresses training with limited and imbalanced data, and \autoref{sec:evaluation} presents the evaluation methodology.


\section{Micro-Expressions}
\label{sec:phenomenon}

Facial expressions are commonly divided into macro-expressions and micro-expressions. Macro-expressions are the ordinary expressions of everyday interaction; they typically last between 0.5 and 4 seconds and involve a large area of the face~\cite{ben2021,xie2022,adegun2020}. Micro-expressions are brief, involuntary and local facial movements that tend to occur when a person attempts to conceal or suppress an emotion, particularly in high-stakes situations~\cite{yan2014,ben2021,xie2022,li2022survey}. Ekman and Friesen described such movements as a form of nonverbal leakage~\cite{ekman1969}, and they have since been studied empirically in concealed and falsified emotional expressions~\cite{porter2008}. Their involuntary character distinguishes them from expressions deliberately posed for a camera~\cite{qu2016}. Comprehensive reviews of automatic micro-expression analysis are provided by Oh et al.~\cite{oh2018}, Ben et al.~\cite{ben2021}, Li et al.~\cite{li2022survey} and Xie et al.~\cite{xie2022}.

\autoref{fig:casme2-sequence} shows an example from CASME~II: a lowering of the inner brow (AU~4, together with AU~9) labelled as disgust, whose apex occurs approximately 110~ms after onset~\cite{yan2014}. The change between frames is barely visible at full-face scale, which illustrates why the phenomenon is difficult to perceive. An observed movement is nevertheless neither proof of deception nor an unambiguous indicator of emotion; the recognition task depends on the annotations of the dataset.

\begin{figure}[htbp]
	\centering
	\includegraphics[width=\textwidth]{fig2_1_casme2_sequence.png}
	\caption[Frame sequence of a CASME II micro-expression]{Frame sequence of a micro-expression from CASME~II, with time in milliseconds from onset. The action units are AU~4 and AU~9, labelled as disgust; the apex occurs at approximately 110~ms. The insets magnify the right inner brow (AU~4). Reproduced from Yan et al.~\cite{yan2014} (Figure~2) under the Creative Commons Attribution (CC~BY) licence.}
	\label{fig:casme2-sequence}
\end{figure}

\subsection{Duration}
\label{sec:duration}

Definitions of micro-expression duration vary. Yan et al.~\cite{yan2014} describe half a second as the generally accepted upper limit, whereas other sources report typical durations of 1/25 to 1/3 of a second~\cite{li2022survey} or 0.065 to 0.5 seconds~\cite{ben2021}. The limit is a working definition rather than a duration filter applied by this pipeline. A brief event yields few observations at ordinary video frame rates, so capture frequency and temporal sampling determine the evidence available for recognition.

\subsection{Intensity}
\label{sec:low-intensity}

Suppressed facial movements may be difficult to perceive even when sufficient frames are recorded~\cite{yan2014}. Because micro-expressions are local~\cite{ben2021}, a small contraction may affect only the region around the brow or mouth while most of the image remains unchanged. Duration limits the number of observations, whereas low intensity limits how clearly the movement can be separated from noise within them. Both properties motivate motion magnification and explicit motion features.

\subsection{Temporal Phases: Onset, Apex and Offset}
\label{sec:onset-apex-offset}

The \textbf{onset} marks the beginning of the movement, the \textbf{apex} its point of greatest annotated intensity, and the \textbf{offset} its end (\autoref{fig:casme2-sequence}). Apex-frame methods use the peak frame as a compact description~\cite{liY2018}. This pipeline instead uses the annotated onset and offset to bound a sequence. The apex annotation is neither read by the model nor used to select frames, so temporal features must be learned from the sampled sequence rather than supplied through a selected peak.

\subsection{Action Units and the Prediction Target}
\label{sec:action-units}

The Facial Action Coding System (FACS)~\cite{ekman1978} describes facial movements through Action Units (AUs), each corresponding to an individual component of muscle movement~\cite{li2022survey}, rather than through emotion categories. CASME~II assigns emotion labels using coded movements together with other annotation evidence~\cite{yan2014}, and the mapping is not one-to-one. This thesis predicts the resulting emotion category only. AU metadata is retained in the coding table but is neither an input nor a training target, and the system does not perform AU recognition.


\section{Data Acquisition and Preprocessing}
\label{sec:recording-to-input}

Input preparation determines which recorded information reaches every subsequent component. In particular, image resizing must be distinguished from face alignment and from temporal interpolation.

\subsection{High-Speed Recording}
\label{sec:high-speed-capture}

The number of frame intervals spanned by a facial movement is approximately
\begin{equation}
	n \approx D f_s,
	\label{eq:frame-intervals}
\end{equation}
where $n$ is the number of frame intervals, $D$ is the duration of the movement in seconds and $f_s$ is the recording rate in frames per second (fps). For an event lasting $D = 0.3$~s, this gives roughly eight to nine intervals at 25--30~fps, against sixty at 200~fps. CASME~II is recorded at 200~fps~\cite{yan2014}. The pipeline takes this rate from its configuration rather than reading it from each clip. After temporal subsampling, the physical spacing between selected frames need not match the nominal rate.

\subsection{Source Frames}
\label{sec:corpus-supplies}

CASME~II provides both the original recorded frames and a registered, cropped release. The study manifest identifies folders of the original $640\times480$ frames rather than the normalised release. The preprocessing performs no face detection, landmark localisation, registration, alignment or cropping of its own. The input therefore includes the full recorded field of view, and neither facial position nor scale is geometrically standardised. The registration described by Yan et al.~\cite{yan2014} must not be attributed to this pipeline.

\subsection{Frame Preparation}
\label{sec:frame-preparation}

Frames are loaded in greyscale and resized by bilinear interpolation to $224\times224$ pixels. Resizing the whole rectangular frame to a square changes its aspect ratio; the result is not a cropped face region. The loader selects the frames between the annotated onset and offset and then samples a fixed-length sequence. No apex localisation is performed.

\subsection{Clip Selection}
\label{sec:clip-selection}

Preprocessing selects CASME~II entries that are coded as micro-expressions, marked as having frames, and have an inclusive onset-to-offset length greater than two frames. The loader must still find readable frames, since the metadata flag alone does not guarantee successful extraction. The recognition dataset additionally requires a matching saved tensor and a label in the selected three-class mapping; the residual \texttt{Others} category is excluded. Preprocessing eligibility and final classification membership are therefore separate filters.


\section{Eulerian Video Magnification}
\label{sec:motion-magnification}

\subsection{The Eulerian Approach}
\label{sec:eulerian-idea}

Eulerian video magnification (EVM), introduced by Wu et al.~\cite{wu2012}, examines intensity changes at fixed image locations rather than explicitly tracking moving points. At each pixel, it isolates temporal variation within a selected frequency band, scales that variation and adds it back to the signal. Bai et al.~\cite{bai2021} describe the amplitude-based mechanism through the small-motion approximation
\begin{equation}
	I(x,t)=f(x+\delta(t))
	\approx f(x)+\delta(t)\frac{\partial f}{\partial x},
	\label{eq:evm-small-motion}
\end{equation}
where $I(x,t)$ is the image intensity at position $x$ and time $t$, $f$ is the intensity profile of the image at the initial time, and $\delta(t)$ is the displacement of that profile at time $t$. Adding $\alpha$ times the motion-related variation gives approximately
\begin{equation}
	\widetilde I(x,t)\approx
	f(x)+(1+\alpha)\delta(t)\frac{\partial f}{\partial x},
	\label{eq:evm-magnified}
\end{equation}
where $\widetilde I(x,t)$ is the magnified intensity and $\alpha$ is the amplification factor. For sufficiently small displacements, intensity amplification therefore resembles a displacement enlarged by the factor $(1+\alpha)$. The approximation does not imply that arbitrary intensity changes represent facial movement, nor that larger amplification always improves recognition.

\subsection{Spatial Decomposition with the Laplacian Pyramid}
\label{sec:laplacian-pyramid}

A Laplacian pyramid~\cite{burt1983} separates an image into spatial scales. At each level, the image is blurred and downsampled, and an upsampled reconstruction is subtracted from the current image to retain the detail at that scale. Repeating this procedure produces a set of Laplacian detail bands and a final coarse Gaussian residual. Reconstruction starts from the residual, upsamples it and adds the detail bands successively. The implementation uses the separable blur kernel $[1,4,6,4,1]/16$ and factor-two downsampling.

\subsection{Temporal Filtering}
\label{sec:temporal-filtering}

Each pixel in a spatial band supplies a time series. A band-pass filter retains the variation within the selected temporal-frequency interval, and the retained signal is amplified before reconstruction. Bai et al.~\cite{bai2021} describe a Butterworth filter for amplitude-based magnification. The present implementation instead sets to zero the Fourier coefficients outside the selected band and then applies the inverse transform.

\subsection{The Amplification Trade-Off}
\label{sec:amplification-tradeoff}

The filter distinguishes frequencies, not causes. Head movement, illumination variation and noise within the pass-band may be amplified together with the facial contraction of interest. A larger $\alpha$ can also violate the small-displacement approximation of \eqref{eq:evm-small-motion} and produce artefacts. Magnification therefore changes the signal presented to flow estimation without guaranteeing a cleaner or more discriminative representation.

\subsection{Implementation Details}
\label{sec:evm-departures}

\autoref{tab:ch2-evm-departures} summarises the implementation choices that affect the interpretation of the magnified input.

\begin{table}[htbp]
    \centering
    \small
    \begin{tabular}{@{}p{0.26\linewidth} p{0.68\linewidth}@{}}
        \hline
        \textbf{Property} & \textbf{Implemented behaviour and consequence} \\
        \hline
        Temporal filter & Hard Fourier mask rather than the Butterworth filter described by Bai et al.~\cite{bai2021}; the pass-band has no smooth roll-off. \\
        Spatial amplification & The same $\alpha$ is used for every Laplacian band, without spatial-frequency damping. The Gaussian residual is neither filtered nor amplified. \\
        Colour treatment & Processing operates on greyscale intensity, with no separate chromatic attenuation. \\
        Pipeline ordering & Each clip is sampled to 33 frames before magnification, but the filter still assumes the nominal 200 fps. Its physical pass-band is therefore clip-dependent. \\
        \hline
    \end{tabular}
    \caption[EVM implementation departures]{Implementation choices that distinguish the magnifier and its position in the pipeline from the reference formulation. These choices are shared by all magnification-enabled configurations.}
    \label{tab:ch2-evm-departures}
\end{table}

The Laplacian pyramid itself is retained. After reconstruction, values are clipped to $[0,255]$ and converted to 8-bit greyscale before motion extraction. Because subsampling changes the frame spacing, the nominal band does not correspond to the same physical-frequency interval for every clip.


\section{Motion Representation: Optical Flow and Optical Strain}
\label{sec:video-to-motion}

\subsection{Motivation for Motion-Based Input}
\label{sec:why-motion}

Optical flow makes displacement explicit instead of requiring the network to infer it from appearance alone. This can reduce reliance on raw image intensity but does not guarantee invariance to identity or illumination: facial geometry, registration errors and brightness changes still affect the estimated field. The resulting representation therefore remains conditioned on the capture and preprocessing procedures.

\subsection{Optical Flow and Brightness Constancy}
\label{sec:optical-flow}

Optical flow estimates the apparent displacement of image content between frames~\cite{horn1981}. It rests on the brightness-constancy assumption that a moving point retains its intensity, as described for micro-expression recognition in OFF-ApexNet~\cite{liong2019a}:
\begin{equation}
	I(x,y,t)=I(x+u\Delta t,y+v\Delta t,t+\Delta t),
	\label{eq:brightness-constancy}
\end{equation}
where $I(x,y,t)$ is the intensity at pixel $(x,y)$ and time $t$, $u$ and $v$ are the horizontal and vertical velocity components, and $\Delta t$ is the time interval between frames. A first-order Taylor expansion gives the optical flow constraint equation~\cite{horn1981}
\begin{equation}
	I_xu+I_yv+I_t=0,
	\label{eq:flow-constraint}
\end{equation}
where $I_x$, $I_y$ and $I_t$ are the partial derivatives of the intensity with respect to $x$, $y$ and $t$. This is a single equation in two unknowns. Under the \textbf{aperture problem}, a local edge constrains the motion normal to the edge, along the intensity gradient, but does not determine the motion parallel to it. Estimators therefore introduce neighbourhood assumptions or regularisation. In the discrete pipeline, $u$ and $v$ denote displacement in pixels between selected frames rather than physical velocity per second. Illumination changes and large inter-frame displacements can violate the underlying approximation.

\subsection{Optical Strain}
\label{sec:optical-strain}

Optical strain describes how neighbouring displacements differ. For the displacement field $\mathbf{u}=[u,v]^{\mathsf T}$, the small-strain approximation used for optical strain~\cite{shreve2011} is
\begin{equation}
	\varepsilon=\tfrac12\left(\nabla\mathbf{u}
	+(\nabla\mathbf{u})^{\mathsf T}\right),
	\qquad
	\varepsilon_{xx}=\frac{\partial u}{\partial x},\quad
	\varepsilon_{yy}=\frac{\partial v}{\partial y},
	\label{eq:strain-tensor}
\end{equation}
\begin{equation}
	\varepsilon_{xy}=\varepsilon_{yx}
	=\tfrac12\left(\frac{\partial u}{\partial y}
	+\frac{\partial v}{\partial x}\right),
	\label{eq:strain-shear}
\end{equation}
where $\varepsilon$ is the $2\times2$ strain tensor, $\nabla\mathbf{u}$ is the spatial gradient (Jacobian) of the displacement field, $\varepsilon_{xx}$ and $\varepsilon_{yy}$ are the normal strain components in the horizontal and vertical directions, and $\varepsilon_{xy}=\varepsilon_{yx}$ is the shear strain component. The normal components describe local stretching or compression, and the off-diagonal terms describe shear. A spatially uniform translation has zero gradient and therefore zero strain. This property does not imply immunity to arbitrary head rotation, projection effects or flow-estimation errors, and differentiation may amplify noise.

The implemented scalar strain magnitude is
\begin{equation}
	\varepsilon_{\mathrm{mag}}
	=\sqrt{\varepsilon_{xx}^2+\varepsilon_{yy}^2
		+\varepsilon_{xy}^2},
	\label{eq:strain-magnitude}
\end{equation}
where $\varepsilon_{\mathrm{mag}}$ is the strain magnitude at a pixel and the remaining symbols are as defined in \eqref{eq:strain-tensor} and \eqref{eq:strain-shear}. This form counts shear once. The Frobenius norm of the symmetric tensor counts both off-diagonal entries, giving $2\varepsilon_{xy}^2$ inside the root. The two forms differ in the relative weighting of shear, not merely by a removable global scale. The scalar discards the direction of deformation but accompanies the original displacement components as a third channel.

\subsection{Construction of the Three-Channel Input Tensor}
\label{sec:three-channel-tensor}

Uniform index sampling selects 33 existing frames from the onset-to-offset sequence. No Temporal Interpolation Model (TIM)~\cite{pfister2011} or other frame synthesis is performed, and repeated indices can occur for short clips. After optional magnification, Farneb\"ack optical flow~\cite{farneback2003} is computed for the 32 adjacent frame pairs, following the estimator choice reported by Zhao et al.~\cite{zhaoS2021}. The OpenCV settings are pyramid scale 0.5, three levels, window size 15, three iterations, polynomial neighbourhood 5, polynomial smoothing 1.2 and flags 0.

The spatial derivatives required for strain are computed with NumPy's \texttt{gradient} function, using unit pixel spacing, central differences in the interior and first-order one-sided differences at the boundaries. Temporal derivatives are not used, and no dedicated strain post-filter is applied. Stacking $[u,v,\varepsilon_{\mathrm{mag}}]$ gives a tensor of shape $(3,32,224,224)$ over the resized full frame. The network receives this tensor rather than the original intensity sequence.

Two normalisations are applied. During extraction, each channel is mapped independently to $[0,1]$ using its extrema across the clip. During loading, each channel is standardised over its $T\times H\times W$ values, after any training augmentation. Small constants in the denominators protect nearly constant channels, which therefore need not have unit variance. Both normalisations operate per sample and use no statistics pooled across training and held-out subjects.


\section{Neural Network Fundamentals}
\label{sec:nn-fundamentals}

A neural network composes parameterised operations to map the motion tensor to emotion-class scores~\cite{lecun2015,goodfellow2016}. This section introduces only the concepts required to interpret the implemented blocks and their controls.

\subsection{Units, Layers and the Forward Pass}
\label{sec:units-layers}

A unit computes
\begin{equation}
	a=\phi(w^{\mathsf T}x+b),
	\label{eq:unit}
\end{equation}
where $x$ is the input vector, $w$ is the learned weight vector, $b$ is the learned bias, $\phi$ is an activation function and $a$ is the resulting activation~\cite{goodfellow2016}. A layer applies several such operations in parallel, and composing layers gives a network. Depth refers to the number of layers and width to their feature dimension. A forward pass evaluates these operations with the current parameters and produces class scores, or logits, from which class probabilities are obtained.

\subsection{Non-Linear Activation Functions}
\label{sec:non-linearity}

Without non-linear operations, a stack of affine layers reduces to a single affine map. Non-linearity allows the composition to represent more complex relationships. The convolutional streams use the rectified linear unit (ReLU)~\cite{nair2010}, and the transformer feed-forward layers use the Gaussian error linear unit (GELU)~\cite{hendrycks2016}:
\begin{equation}
	\operatorname{ReLU}(x)=\max(0,x),
	\qquad
	\operatorname{GELU}(x)=x\,\Phi(x),
	\label{eq:activations}
\end{equation}
where $x$ is the input to the activation and $\Phi$ is the cumulative distribution function of the standard normal distribution. Both are fixed functions, although they act on features learned by the surrounding layers.

\subsection{Gradient-Based Learning}
\label{sec:gradient-descent}

A loss function compares predictions with labels. Backpropagation~\cite{rumelhart1986} applies the chain rule through the network to compute the derivatives of the loss with respect to its parameters. An optimiser uses these gradients to update the parameters, with the learning rate controlling the size of each update. Mini-batches provide successive updates, and an epoch completes the configured number of batch draws. When sampling is performed with replacement, an epoch need not visit every distinct clip.

During evaluation, the parameters are held fixed and training-only stochastic behaviour such as dropout is disabled. A low training loss alone does not establish useful learning, because the relevant objective is prediction on unseen subjects. A validation set used to select an epoch therefore has a different role from a test set reserved for final assessment.

\subsection{Convolution}
\label{sec:convolution}

A convolution applies the same learned kernel at multiple spatial positions~\cite{lecun1998}. This weight sharing allows local patterns to be detected without learning a separate detector for every image location. Each output channel is a feature map, and each kernel combines all input channels over its local window. For a two-dimensional layer with one bias per output channel, the number of parameters is
\begin{equation}
	N_{\mathrm{params}}=C_{\mathrm{out}}
	(C_{\mathrm{in}}k_Hk_W+1),
	\label{eq:conv-params}
\end{equation}
where $C_{\mathrm{in}}$ and $C_{\mathrm{out}}$ are the numbers of input and output channels, $k_H$ and $k_W$ are the kernel height and width, and the added one accounts for the bias of each output channel. The parameter count is independent of the image dimensions, but computation and activation memory increase with the number of positions processed. The stride sets the spacing between kernel applications, and padding controls boundary treatment and output size. Stacking layers enlarges the receptive field, that is, the input region that can influence a later activation.

\subsection{Pooling}
\label{sec:pooling}

Pooling replaces a neighbourhood with a fixed summary such as its maximum or mean~\cite{lecun1998,goodfellow2016}. A $2\times2$ max-pooling operation with stride two halves the spatial dimensions and retains the strongest response in each window. Adaptive average pooling instead produces a specified output grid. Neither operation learns weights. Both discard detail while reducing the computation of subsequent layers; any later learned projection is a separate operation with its own parameters.

\subsection{Model Capacity and Small Datasets}
\label{sec:capacity}

Capacity describes the range of functions a model can represent~\cite{goodfellow2016}. The parameter count is one indicator, but architecture, constraints and training procedure also matter. A model can overfit subject-specific appearance, noise or recording conditions, achieving a low training loss without generalising. Having more parameters than examples does not by itself prove overfitting, just as few parameters do not guarantee good predictions.

The small size of the CASME~II training pool motivates controlled architectural comparisons and regularisation. Their effectiveness must be assessed on held-out subjects, with checkpoint selection kept separate from final testing where possible. Because the analytical flow and strain transformations already expose motion structure, the learned network does not start from raw pixels.


\section{Network Building Blocks}
\label{sec:building-blocks}

The model combines a spatial stem, optional activation re-weighting, an optional temporal encoder and a common classification head. The stem consists of three streams with separate learned weights, one each for horizontal flow, vertical flow and strain. Disabling a component replaces a specified operation while preserving compatible input and output shapes.

\subsection{Three-Dimensional Convolution and Kernel Shape}
\label{sec:conv3d}

A three-dimensional convolution~\cite{ji2013,tran2015} operates on tensors of shape $(B,C,T,H,W)$, denoting batch size, channels, time, height and width. A kernel of shape $(k_T,k_H,k_W)$ mixes information across time only when $k_T>1$. The implemented streams use $(1,3,3)$ kernels with $(0,1,1)$ padding, followed by $(1,2,2)$ max-pooling. The convolutions therefore preserve the temporal length and learn spatial filters, while pooling halves the height and width.

The use of \texttt{Conv3d} does not establish that temporal convolution was tested. The complete stream can nevertheless depend on multiple frames, because batch normalisation during training and SimAM, whenever enabled, use statistics that span time. This distinction separates temporal convolution from other forms of whole-clip dependence.

\subsection{Normalisation and Regularisation}
\label{sec:normalisation}

Batch normalisation~\cite{ioffe2015} rescales each channel using a mean and variance, followed by a learned scale and shift:
\begin{equation}
	\widehat x=\frac{x-\mu_c}{\sqrt{\sigma_c^2+\epsilon}},
	\qquad y=\gamma_c\widehat x+\beta_c,
	\label{eq:batchnorm}
\end{equation}
where $x$ is an activation in channel $c$, $\mu_c$ and $\sigma_c^2$ are the mean and variance of that channel, $\epsilon$ is a small constant for numerical stability, $\widehat x$ is the normalised activation, $\gamma_c$ and $\beta_c$ are the learned scale and shift of the channel, and $y$ is the output. \texttt{BatchNorm3d} computes the training statistics over the $(B,T,H,W)$ dimensions of each channel, and evaluation uses running estimates accumulated during training. These activation statistics differ from the per-sample input standardisation performed by the data loader.

Dropout~\cite{srivastava2014} randomly deactivates units during training to reduce overfitting. \texttt{Dropout3d} removes entire feature channels for a sample~\cite{tompson2015} rather than masking individual spatial values. The retained channels are rescaled, and evaluation disables the random mask. Layer normalisation~\cite{ba2016} instead standardises the feature dimension of each token independently of other samples. It appears inside the transformer and in the classification head, where it is followed by dropout and a linear projection to class logits.

\subsection{Parameter-Free Attention: SimAM}
\label{sec:simam}

SimAM assigns larger weights to activations that differ from their channel context~\cite{yang2021}. Following the reference code of the original paper, the statistics of each context are
\begin{equation}
	\mu=\frac1M\sum_i x_i,\qquad
	v=\frac1{M-1}\sum_i(x_i-\mu)^2,
	\label{eq:simam-stats}
\end{equation}
and each activation is re-weighted as
\begin{equation}
	E_i^{-1}=\frac{(x_i-\mu)^2}{4(v+\lambda)}+\frac12,
	\qquad \widetilde x_i=x_i\operatorname{sigmoid}(E_i^{-1}),
	\label{eq:simam-gate}
\end{equation}
where $x_i$ is the activation at position $i$ within a channel, $M$ is the number of positions in that channel, $\mu$ and $v$ are the mean and variance of the channel, $\lambda$ is a regularisation constant, $E_i^{-1}$ is the reciprocal of the minimal energy of position $i$, and $\widetilde x_i$ is the re-weighted activation. The reciprocal energy measures distinctiveness, and the sigmoid bounds the resulting multiplicative weight while preserving the order of the scores. The printed shared-variance definition in the paper uses $M$, whereas the code in its Figure~3 uses $M-1$; the implementation follows the latter.

In this model, $M=THW$ within each channel, rather than the $HW$ of the original image formulation~\cite{yang2021}, and $\lambda=10^{-4}$. SimAM acts on each convolutional stream before concatenation. It adds no learned parameters but still performs reductions and element-wise computations. Whole-clip statistics do not guarantee that apex frames receive the highest weights, since unusual activation values may also reflect noise.

\subsection{Self-Attention and the Transformer Encoder}
\label{sec:self-attention}

In self-attention, each input token is projected into a query, a key and a value. Scaled dot-product attention~\cite{vaswani2017}, in the form also presented by Dosovitskiy et al.~\cite{dosovitskiy2021}, computes
\begin{equation}
	\operatorname{Attention}(Q,K,V)
	=\operatorname{softmax}\left(\frac{QK^{\mathsf T}}{\sqrt{d_k}}\right)V,
	\label{eq:attention}
\end{equation}
where $Q$, $K$ and $V$ are the matrices of queries, keys and values, with one row per token, $d_k$ is the dimension of the key vectors, and the softmax is applied to each row. The queries and keys determine the contribution of each value to each output. The scaling by $\sqrt{d_k}$ moderates the magnitude of the dot products as the key dimension grows. Multiple attention heads learn projections into different subspaces, whose outputs are concatenated and linearly combined~\cite{vaswani2017}.

Without position information, self-attention is permutation-equivariant: permuting the input tokens permutes their outputs correspondingly, so the temporal order is not identified. Fixed sinusoidal position vectors~\cite{vaswani2017} provide this information in the present model. Although the sinusoidal functions can be evaluated at other positions, generalisation to unseen sequence lengths is not guaranteed, and the implementation uses a finite precomputed buffer.

The temporal encoder alternates attention and position-wise feed-forward sublayers with residual connections~\cite{he2016}. It uses pre-norm placement, applying layer normalisation before each sublayer~\cite{xiong2020}. Its output is averaged over time before classification. This mean summarises temporal features that are already contextualised, and therefore differs from averaging the input tokens without an encoder.

\subsection{Replacement Operations for Disabled Components}
\label{sec:fallbacks}

\autoref{tab:ch2-component-controls} defines the operation that replaces each component when it is disabled.

\begin{table}[htbp]
    \centering
    \small
    \begin{tabular}{@{}p{0.22\linewidth} p{0.35\linewidth} p{0.35\linewidth}@{}}
        \hline
        \textbf{Component} & \textbf{Enabled} & \textbf{Disabled} \\
        \hline
        Spatial stem & Three learned convolutional streams & Per-frame $4\times4$ average pooling and a learned $48\rightarrow96$ linear projection \\
        SimAM & Per-stream activation re-weighting & Unmodified stream outputs; the model enables SimAM only together with its CNN stem \\
        Temporal encoder & Position encoding, transformer layers, then temporal mean & Direct mean of the input sequence \\
        \hline
    \end{tabular}
    \caption[Enabled components and their controls]{Operations compared by the model switches. The classification head is shared in structure; compatible feature dimensions do not imply equal parameter counts.}
    \label{tab:ch2-component-controls}
\end{table}

Pooling and direct temporal averaging have no learned weights. The projection used without the CNN stem does: including its bias, it contains $48\times96+96=4{,}704$ parameters. The replacements preserve the interfaces required by the experiment while changing the operations and the capacity. They do not isolate architecture independently of parameter count or computational cost.


\section{Training with Limited and Imbalanced Data}
\label{sec:scarcity-skew}

Micro-expression datasets are small and their emotion classes are unevenly represented. This section introduces the loss functions, sampling strategy, augmentation and optimisation methods used to train under these conditions.

\subsection{Cross-Entropy under Class Imbalance}
\label{sec:skew-cross-entropy}

The cross-entropy loss of a single example is
\begin{equation}
	L_{\mathrm{CE}}=-\log p_t,
	\label{eq:cross-entropy}
\end{equation}
where $p_t$ is the probability predicted for the true class $t$. Under natural sampling, a frequent class contributes more terms to the training objective, which biases learning towards the majority class~\cite{he2009}. A low overall loss therefore does not imply that rare classes are recognised well.

\subsection{Focal Loss}
\label{sec:focal-loss}

Focal loss, introduced by Lin et al.~\cite{lin2017} and applied to micro-expression recognition by Zhao et al.~\cite{zhaoS2021}, is defined as
\begin{equation}
	\operatorname{FL}(p_t)=-\alpha_t(1-p_t)^\gamma\log p_t,
	\label{eq:focal-loss}
\end{equation}
where $p_t$ is the predicted probability of the true class, $\gamma\geq0$ is the focusing parameter and $\alpha_t$ is an optional weighting factor for the true class. For $\gamma>0$, the modulating factor $(1-p_t)^\gamma$ reduces the contribution of confidently classified examples relative to difficult ones. This factor depends on prediction difficulty, not on class frequency; any frequency correction must come from $\alpha_t$ or from the sampling strategy. This thesis uses $\gamma=2$, the value reported by Zhao et al.~\cite{zhaoS2021}.

\subsection{Class-Balanced Sampling}
\label{sec:balanced-sampling}
\label{sec:combining-both}

Oversampling the minority classes is a common data-level response to class imbalance~\cite{he2009,buda2018}. Each training example is assigned the sampling weight
\begin{equation}
	w_i=\frac{1}{n_{c_i}},
	\label{eq:sampler-weight}
\end{equation}
where $w_i$ is the weight of training example $i$, $c_i$ is its class, and $n_{c_i}$ is the number of training examples of that class. Drawing with replacement according to these weights makes each class equally likely in expectation. Balanced sampling changes how often each class is seen, not the number of distinct examples, and it is applied to training only; evaluation uses the natural class distribution of the held-out data.

Class weights in the loss ($\alpha_t$ in \eqref{eq:focal-loss}) and balanced sampling correct the same imbalance at different points. Applying both corrects it twice. In this thesis the frequency correction is therefore assigned to the sampler alone, and the loss carries no class weights.

\subsection{Label Smoothing}
\label{sec:label-smoothing}

Label smoothing~\cite{szegedy2016} replaces the one-hot target with a mixture of the one-hot target and a uniform distribution. For $C$ classes, the smoothed target is
\begin{equation}
	q_c=(1-\epsilon)\mathbf{1}[c=t]+\frac{\epsilon}{C},
	\label{eq:smoothed-target}
\end{equation}
and the corresponding cross-entropy term is
\begin{equation}
	L_{\mathrm{smooth}}
	=-(1-\epsilon)\log p_t
	-\frac{\epsilon}{C}\sum_{c=1}^{C}\log p_c,
	\label{eq:label-smoothing}
\end{equation}
where $q_c$ is the target probability of class $c$, $\epsilon$ is the smoothing strength, $t$ is the true class, $\mathbf{1}[c=t]$ equals one when $c=t$ and zero otherwise, and $p_c$ is the predicted probability of class $c$. Label smoothing discourages over-confident predictions~\cite{szegedy2016}. In this thesis, the focal factor $(1-p_t)^\gamma$ of \eqref{eq:focal-loss} multiplies the smoothed term \eqref{eq:label-smoothing}, with $p_t$ taken from the unsmoothed prediction.

\subsection{Data Augmentation}
\label{sec:augmentation}

Data augmentation enlarges the effective training set through label-preserving transformations, and Xia et al.~\cite{xia2020a} motivate it for limited and imbalanced micro-expression data. The augmentation used here is a random horizontal flip, applied with probability one half during training only. Because mirroring a clip reverses the direction of horizontal motion, the horizontal flow channel $u$ is negated whenever the clip is flipped.

\subsection{Optimisation}
\label{sec:optimisation}

AdamW~\cite{loshchilov2019}, a variant of Adam~\cite{kingma2015}, updates the parameters with weight decay decoupled from the gradient-based update. The learning rate is increased linearly during an initial warmup phase and then reduced by cosine annealing~\cite{loshchilov2017}, a schedule also used for micro-expression recognition by Zhang et al.~\cite{zhang2022}. Mixed-precision training~\cite{micikevicius2018} performs most computations in 16-bit floating point to reduce memory use, with gradient scaling to prevent small gradient values from underflowing.


\section{Evaluation Methodology}
\label{sec:evaluation}

This section defines the metrics and the cross-validation protocol used to evaluate the models.

\subsection{Confusion Matrix, Precision, Recall and F1}
\label{sec:confusion-matrix}

In a confusion matrix, entry $(i,j)$ counts the examples whose true class is $i$ and whose predicted class is $j$, so the diagonal contains the correct predictions. The class-wise metrics follow OFF-ApexNet~\cite{liong2019a} and the MEGC 2019 protocol~\cite{see2019}:
\begin{equation}
	P_c=\frac{TP_c}{TP_c+FP_c},\qquad
	R_c=\frac{TP_c}{TP_c+FN_c},\qquad
	F1_c=\frac{2TP_c}{2TP_c+FP_c+FN_c},
	\label{eq:prf}
\end{equation}
where $P_c$, $R_c$ and $F1_c$ are the precision, recall and F1 score of class $c$; $TP_c$ (true positives) counts correct predictions of $c$; $FP_c$ (false positives) counts incorrect predictions of $c$; and $FN_c$ (false negatives) counts examples of $c$ predicted as another class. Precision penalises false alarms and recall penalises misses. F1 is their harmonic mean, so a high value requires both to be high. When a denominator is zero, the score is set to zero.

\subsection{Macro and Micro Averaging}
\label{sec:macro-micro}

Macro F1 gives each class equal weight, using the mean of per-class F1 scores specified by MEGC 2019~\cite{see2019}:
\begin{equation}
	F1_{\mathrm{macro}}=\frac1C\sum_{c=1}^{C}F1_c,
	\label{eq:macro-f1}
\end{equation}
where $C$ is the number of classes and $F1_c$ is defined in \eqref{eq:prf}. Micro F1 instead sums the true positives, false positives and false negatives across classes before computing the score~\cite{sokolova2009}. In single-label multiclass classification, micro F1 equals accuracy and can therefore be dominated by a frequent class. Macro F1 is the more informative measure under class imbalance.

\subsection{Pooled and Per-Fold Averaging}
\label{sec:pooled-vs-fold}

Under cross-validation, macro F1 can be computed in two ways: by pooling the confusion matrices of all folds before computing the class F1 scores, or by computing macro F1 in each fold and averaging the fold scores. Because F1 is non-linear, the two need not agree. If a fold contains no examples of a class, that class contributes zero to the fold score even when every example is classified correctly, so averaging fold scores mixes classifier behaviour with fold composition. Pooling avoids this. The Unweighted F1 (UF1) of MEGC 2019 is defined through the pooled counts~\cite{see2019}, and pooled macro F1 is the primary metric in this thesis.

\subsection{Cross-Validation and Subject Independence}
\label{sec:cross-validation}

Cross-validation repeatedly fits a model to a training partition and evaluates it on held-out examples. Leave-one-subject-out (LOSO) cross-validation holds out all clips of one subject in each fold and trains on the remaining subjects. It is the native protocol of CASME~II~\cite{yan2014} and the protocol mandated by MEGC 2019~\cite{see2019}. Because no subject appears in both partitions, the model cannot exploit familiarity with a subject's identity or recording conditions.

Two further principles govern how such results can be interpreted. First, if the held-out data are also used for model selection, such as choosing the best training epoch, the resulting estimate is optimistically biased; an inner validation split drawn from the training subjects avoids this~\cite{varma2006}. Second, testing whether two classifiers evaluated on the same examples differ significantly requires a paired test such as McNemar's test~\cite{mcnemar1947}, which needs the prediction for every individual example.


\section{Summary}
\label{sec:ch2-conclusion}

This chapter introduced the concepts on which the thesis builds. It defined micro-expressions and the properties that make them difficult to recognise: short duration and low intensity, and the relationship between action units and emotion labels. It described the CASME~II recordings and their preparation, Eulerian video magnification, and the optical flow and optical strain representations that form the model input. It then introduced the neural network concepts and building blocks used by the model, the methods used to train it on limited and imbalanced data, and the metrics and leave-one-subject-out protocol used to evaluate it.

\autoref{ch:literature} reviews the published work on each of these components. For each stage of the pipeline, it examines the evidence for its reported contribution to micro-expression recognition and identifies the gap that this thesis addresses.
